\documentclass[11pt,a4paper]{amsart}
\usepackage{cgeometry}
\addbibresource{ComplexGeometry.bib}
\defbibheading{bibliography}[\refname]{\section*{#1}}
\setlength{\bibitemsep}{0pt}
\renewcommand*{\bibfont}{\footnotesize}

\allowdisplaybreaks[2]
\emergencystretch=1em

%%%%%%%%%%%%%%%%%%%%%%%%%%%%%%%%%%%%%%%%%%%%%%%%%%%%%%%%%%%%%%%%%%%%%%%%%%%%%
%% Notation
%%%%%%%%%%%%%%%%%%%%%%%%%%%%%%%%%%%%%%%%%%%%%%%%%%%%%%%%%%%%%%%%%%%%%%%%%%%%%
\newcommand{\EnK}{E_{\mathrm{nK}}}
\newcommand{\Enn}{E_{\mathrm{nn}}}
\newcommand{\trc}{\operatorname{Tr}}            % trace operator
\newcommand{\dS}{d_{S}}                          % the d_S pseudometric
\DeclareMathOperator{\Exc}{Exc}
\DeclareMathOperator{\im}{im}
\newcommand{\rest}[1]{\operatorname{Res}_{#1}}   % restriction of an ideal
\newcommand{\lessing}{\preceq}                   % more singular than
\newcommand{\lessP}{\preceq_{P}}
\newcommand{\lessI}{\preceq_{\mathcal I}}
\newcommand{\simP}{\sim_{P}}
\newcommand{\simI}{\sim_{\mathcal I}}

\theoremstyle{definition}
\newtheorem{setup}[theorem]{Setup}

\title{A transcendental Nakamaye theorem}
\author{Duc Viet Vu \and Mingchen Xia}
\date{}

\begin{document}
\begin{abstract}
Let $X$ be a connected compact K\"ahler manifold and let $\alpha$ be a
pseudoeffective real $(1,1)$-class on $X$. We prove that every
positive-dimensional irreducible component $V$ of the non-K\"ahler locus
$\EnK(\alpha)$ has vanishing numerical restricted volume,
\[
  \langle\alpha^{\dim V}\rangle_{X\mid V}=0.
\]
Consequently, the non-K\"ahler locus of $\alpha$ coincides with the null
locus of $\alpha$ defined through numerical restricted volumes. This
confirms a conjecture of Collins and Tosatti, and gives a transcendental
version of the theorems of Nakamaye and of
Ein--Lazarsfeld--Musta\c{t}\u{a}--Nakamaye--Popa on augmented base loci.
\end{abstract}
\maketitle


\section{Introduction}

Let $X$ be a connected compact K\"ahler manifold of dimension $n$. For a big
class $\alpha\in H^{1,1}(X,\mathbf R)$, Boucksom \cite{Bou04} introduced
the \emph{non-K\"ahler locus}
\[
  \EnK(\alpha)=\bigcap_{T}E_+(T),
\]
where $T$ runs over the K\"ahler currents in $\alpha$ and $E_+(T)$
is the set of points at which the Lelong number of $T$ is positive.
By \cite[Theorem~3.17]{Bou04}, there is a K\"ahler current
$T\in\alpha$ with analytic singularities such that
$E_+(T)=\EnK(\alpha)$; in particular, $\EnK(\alpha)$ is a proper analytic
subset of $X$. When $X$ is projective and $\alpha=c_1(L)$ for a big
line bundle $L$, the non-K\"ahler locus is the augmented base
locus $B_+(L)$. Following \cite{CT22}, we set $\EnK(\alpha)=X$
when $\alpha$ is pseudoeffective but not big.

Nakamaye's theorem \cite{Nak00} describes the augmented base locus of a nef
and big line bundle $L$ on a projective manifold as the \emph{null locus}
of $L$, the union of the positive-dimensional irreducible
subvarieties $V\subset X$ with $(L^{\dim V}\cdot V)=0$. It was extended
to arbitrary big divisors $D$ by
Ein--Lazarsfeld--Musta\c{t}\u{a}--Nakamaye--Popa \cite{ELMNP09}, who showed
that $B_+(D)$ is the union of the positive-dimensional subvarieties along
which the restricted volume $\vol_{X\mid V}(D)$ vanishes. On the
transcendental side, Collins and Tosatti \cite{CT15} proved that
$\EnK(\alpha)=\Null(\alpha)$ for every nef and big class $\alpha$ on a
compact K\"ahler manifold, where
$\Null(\alpha)$ is the union of the positive-dimensional irreducible
subvarieties $V$ with $\int_V\alpha^{\dim V}=0$.

To formulate a statement valid for all pseudoeffective classes,
Collins and Tosatti \cite{CT22} introduced the \emph{numerical restricted
volume}
$\langle\alpha^{\dim V}\rangle_{X\mid V}$ of a pseudoeffective class
$\alpha$ along a positive-dimensional irreducible analytic subvariety
$V\subset X$ (see \cref{sec:restricted-volumes}), defined the null locus
\begin{equation}\label{eq:null-locus}
  \Null(\alpha)
  =\bigcup_{\substack{V\subset X\ \text{irreducible},\ \dim V>0\\
      \langle\alpha^{\dim V}\rangle_{X\mid V}=0}}V,
\end{equation}
and conjectured that $\EnK(\alpha)=\Null(\alpha)$ for every pseudoeffective
class $\alpha$ \cite[Conjecture~1.1]{CT22}. The inclusion
$\Null(\alpha)\subset\EnK(\alpha)$ always holds
(\cref{sec:restricted-volumes}\ref{it:rv-positive}). They proved the conjecture
when $\alpha$ admits a Zariski decomposition, in particular whenever
$\dim X=2$ \cite[Theorem~1.2]{CT22}; when $X$ is projective and
$\alpha=c_1(D)$ for a big $\mathbf R$-divisor $D$, the conjecture reduces
to the theorem of \cite{ELMNP09} \cite[Proposition~1.4]{CT22}. Since the
non-K\"ahler locus of a big class has no isolated points
(\cref{lem:no-isolated-points}), the conjecture follows once the
numerical restricted volume is known to vanish along every irreducible
component of $\EnK(\alpha)$. This is what we prove.

\begin{samepage}
\begin{maintheorem}\label{thm:main}
Let $X$ be a connected compact K\"ahler manifold and let $\alpha$ be a
pseudoeffective real $(1,1)$-class on $X$. If $V$ is a
positive-dimensional irreducible component of $\EnK(\alpha)$, then
\[
  \langle\alpha^{\dim V}\rangle_{X\mid V}=0.
\]
In particular, $\EnK(\alpha)=\Null(\alpha)$.
\end{maintheorem}
\end{samepage}

Collins and Tosatti derived several consequences of their conjecture
\cite[Theorems~1.6, 1.7 and~1.8]{CT22}, concerning the moving Seshadri
constants of big classes, their continuity on the big cone, and the
non-K\"ahler locus of classes containing a positive current which is a
K\"ahler metric on an open set; by \cref{thm:main}, these statements now
hold unconditionally.

\subsection*{Where the difficulty lies}
Let $\alpha$ be big and let $V$ be a smooth irreducible component of
$\EnK(\alpha)$ of dimension $k>0$, not contained in the non-nef locus
$\Enn(\alpha)$; the general case reduces to this one
(\cref{sec:setup}). Assume that
$h\coloneqq\langle\alpha^k\rangle_{X\mid V}>0$. We shall reach a
contradiction by constructing, on a modification $\pi\colon Y\to X$, a
K\"ahler current in $\pi^*\alpha$ whose potential is bounded near some point
of the strict transform of $V$.

In the algebraic setting, the theorem of \cite{ELMNP09} rests on the
extension of sections of line bundles from $V$ to $X$; no such tool is
available for transcendental classes. What remains is the gluing
mechanism already used in the nef case \cite{CT15}: a K\"ahler current
$T_U\in\alpha$ defined near $V$ and a global K\"ahler current
$T_G\in\alpha$ singular along $V$ can be glued, by taking the maximum of
their potentials, as soon as the potential of $T_U$ is bounded above by
that of $T_G$, up to a constant, near $V$ but off $V$
(\cref{prop:gluing}). We call this the \emph{collar condition}. The
natural candidate for $T_U$ is an extension of a K\"ahler current on $V$
in the class $\alpha|_V$. When $\alpha$ is nef, $T_G$ may be chosen with
arbitrarily mild singularities and the collar condition costs nothing.
For a general big class it is the heart of the matter, and three
problems arise.

\begin{enumerate}[label=(\arabic*),leftmargin=*,itemsep=.2em,topsep=.3em]
\item \emph{Singularities are forced on $V$.} Every closed positive
  current in $\alpha$ is at least as singular as the potential $V_\theta$ with minimal
  singularities ($\theta$ a smooth representative of $\alpha$), which may be singular along subvarieties of $V$, since
  $V$ may meet $\Enn(\alpha)$. Hence $T_U$, and the current on $V$ from
  which it is built, must be at least as singular as the restrictions
  $\mathfrak b_q$ to $V$ of the multiplier ideals $\mathcal I(qV_\theta)$.
  The trace operator of \cite{DX24} provides the relevant object: the
  trace $u$ of $V_\theta$ on $V$ has non-pluripolar mass $h$, and its
  multiplier ideals satisfy $\mathcal I(qu)\subset\mathfrak b_q$. But
  $u$ is not a K\"ahler current, and a positive mass only says that
  $\alpha|_V$ is big. A domination lemma (\cref{lem:domination}),
  based on the bigness criterion of Xiao and Popovici
  \cite{Xiao15,Pop16}, turns a lower bound on the mass into a K\"ahler
  current on $V$ that is at least as singular as $u$ in the sense of
  multiplier ideals, with uniform positivity.
\item \emph{Singular currents do not extend.} There is no general
  procedure extending a singular K\"ahler current from $V$ to a
  neighborhood of $V$. We therefore make its singularities divisorial on a
  modification $\pi\colon Y\to X$ obtained by blowing up centers contained
  in, or disjoint from, the strict transforms of $V$. On the strict transform $W$ of $V$ the current
  becomes a K\"ahler class plus an effective divisor; the class extends to a
  K\"ahler form near $W$, and the divisor extends as soon as it is the
  restriction to $W$ of an effective divisor $D$ on $Y$. The collar
  condition then becomes a comparison, near $W$, between $D$ and the
  singularities of $\pi^*T_G$.
\item \emph{Blowing up has a price.} To be a K\"ahler current, $T_G$ is
  built from a regularization of $V_\theta$ at a level $M$, which loses
  positivity of order $1/M$, corrected by a small multiple $(L/M)R$ of a
  K\"ahler current $R$ singular along $\EnK(\alpha)\supset V$. On $Y$ we
  have $\mathcal I_V\mathcal O_Y=\mathcal O_Y(-H)\cdot\mathcal I_W$ for an
  exceptional divisor $H$ which meets $W$, so $\pi^*T_G$ is singular along
  $H$, with a coefficient of order $1/M$. The collar condition forces the
  current on $W$ to be more singular than the trace by an amount of order
  $1/M$ along $H|_W$, and the trace offers no such margin.
\end{enumerate}

\subsection*{Strategy}
The margin is created by a rooftop envelope at two levels $m\ll M$. Let
$\phi_m$ be the restriction to $V$ of a Demailly approximation of
$V_\theta$ at level $m$; it is less singular than $u$, and its
singularities are those of $\mathfrak b_m$. For an integer $C$ chosen
before $m$ and $M$, the theorem on rooftop envelopes
(\cref{thm:rooftop}) yields a potential
\[
  v\le Cu-(C-1)\phi_m
\]
whose non-pluripolar mass is bounded below independently of $m$, and
hence a K\"ahler current $K$ on $V$ which is at least as singular as $v$.
Writing $M=C\ell$ and $p=(C-1)\ell/m$, the multiplier ideal
$\mathfrak k=\mathcal I(\ell\varphi_K)$ of the potential
$\varphi_K$ of $K$ at level $\ell$ satisfies
(\cref{lem:ideal-comparison})
\[
  \mathfrak k\,\mathfrak b_m^{\,p}\subset\mathcal I(Mu)\subset\mathfrak b_M .
\]
On a modification of $V$ on which $\mathfrak k$, $\mathfrak b_m$ and
$\mathfrak b_M$ become invertible, with divisors $G_K$, $G_m$ and $G_M$,
this means that
\[
  \frac1\ell\,G_K=\frac1M\,G_M+\frac{C-1}M\Bigl(G_M-\frac Mm\,G_m\Bigr)
  +\frac CM\,F
\]
with $F\ge0$ and $G_M-(M/m)G_m\ge0$. The first term matches the
singularities of the regularization of $V_\theta$ at level $M$; the other
two terms are the margin. By a principalization relative to the strict
transform of $V$ (\cref{lem:relative-principalization}) we construct $Y$
so that $H$ is the sum of three exceptional divisors whose restrictions
to $W$ are $G_m$, $G_M-(M/m)G_m$ and $F$. The last two pieces are paid
for by the margin, because $C$ is chosen so that $(C-1)/M$ dominates the
coefficient of $H$ in $\pi^*T_G$. The first piece does not depend on $M$; it is absorbed
by a global quasi-plurisubharmonic function with loss of positivity of
order $1/M$. This is why two levels are needed.

The order in which the constants are chosen is essential: the constants
governing positivity are fixed before $M$, while singularities are
compared at the level $M$. This is possible because Demailly's
regularization loses positivity at the rate $O(1/M)$, with a constant
independent of the function being regularized
(\cref{lem:regularization}): the weight of $R$, and hence the margin
needed, is then of order $1/M$, like the margin produced by the rooftop
construction. The order of choices is summarized in
\cref{rem:order}.

\subsection*{Organization}
\Cref{sec:prelim} collects notation and the results from pluripotential
theory used in the proof; we follow the lecture notes \cite{XiaSGPT},
whose results we use freely. \Cref{sec:aux} contains the auxiliary
results: the bigness criterion of Xiao and Popovici, the domination lemma,
a patching lemma with its two applications (the quantitative
regularization lemma and functions with prescribed singularities), the
relative principalization lemma, a lemma on K\"ahler forms on blowups,
and the gluing proposition. The proof of \cref{thm:main} occupies
\cref{sec:setup,sec:positive-current,sec:extension}:
\cref{sec:setup} reduces the theorem to a statement about the trace,
\cref{sec:positive-current} constructs the K\"ahler current $K$ on $V$,
and \cref{sec:extension} extends it across $V$ and concludes.

\subsection*{Acknowledgments}
The initial proof was due to the
\href{https://github.com/frenzymath/Danus}{Danus agent} (improved by Felix
Ye), using the \texttt{gpt-5.6-sol} model. The authors then simplified
the construction and improved the presentation, and are fully responsible
for all assertions in this paper.

The second-named author is supported by the National Key R\&D Program
of China under Grant No.~2025YFA1018200.

\section{Preliminaries}\label{sec:prelim}

Throughout, $X$ denotes a connected compact K\"ahler manifold of
dimension $n$ and $\omega$ a K\"ahler form on $X$. Submanifolds are
closed; in particular, submanifolds of compact manifolds are compact. We
write $\ddc=(2\pi)^{-1}\mathrm{i}\partial\bar\partial$, so that
$\ddc\log|z|^2$ is the current of integration at the origin of
$\mathbf C$. A \emph{modification} is a finite composition of blowups
with smooth centers between compact complex manifolds. A \emph{prime
divisor over $X$} is a prime divisor $E$ on some modification
$\pi\colon Y\to X$. For a quasi-plurisubharmonic (quasi-psh) function
$\varphi$ on $X$, we write $\nu(\varphi,x)$ for its Lelong number at
$x\in X$, $\nu(\varphi,Z)$ for its generic Lelong number along an
irreducible analytic subset $Z\subset X$, and $\nu(\varphi,E)$ for the
generic Lelong number of $\pi^*\varphi$ along a prime divisor $E$ over
$X$. The same notation is used for closed positive $(1,1)$-currents.
If $D$ is a divisor on a modification and $f$ is a further modification,
$f^*D$ denotes the total transform of $D$.

\subsection{Quasi-psh functions, multiplier ideals and analytic
singularities}\label{sec:qpsh}

For a smooth closed real $(1,1)$-form $\theta$ on $X$, we denote by
$\PSH(X,\theta)$ the set of quasi-psh functions $\varphi$ on $X$ with
$\theta_\varphi\coloneqq \theta+\ddc\varphi\ge0$; equivalently, if
$\theta=\ddc g$ locally, then $g+\varphi$ is psh. For quasi-psh functions
$\varphi,\psi$ we write $\varphi\lessing\psi$, and say that $\varphi$ is
\emph{more singular} than $\psi$, if $\varphi\le\psi+C$ for a constant
$C$; we write $\varphi\sim\psi$ if $\varphi\lessing\psi$ and
$\psi\lessing\varphi$.

For a quasi-psh function $\varphi$ and $\lambda>0$, the
\emph{multiplier ideal} $\mathcal I(\lambda\varphi)\subset\mathcal O_X$
\cite{Nad90} is the sheaf of germs of holomorphic functions $f$ with
$|f|^2\mathrm{e}^{-\lambda\varphi}\in L^1_{\mathrm{loc}}$. It is coherent
and integrally closed, and by the subadditivity theorem \cite{DEL00},
\begin{equation}\label{eq:subadditivity}
  \mathcal I\bigl((\lambda+\lambda')\varphi\bigr)
  \subset\mathcal I(\lambda\varphi)\,\mathcal I(\lambda'\varphi),
  \qquad \lambda,\lambda'>0.
\end{equation}

For a coherent ideal $\mathfrak a$ with local generators
$f_1,\dots,f_N$, we set $\psi_{\mathfrak a}\coloneqq \log\sum_j|f_j|^2$.
We write $O(1)$ for a locally bounded function; a different choice of
generators changes $\psi_{\mathfrak a}$ by $O(1)$, and inequalities
between such locally defined functions are understood locally, with
constants that can be chosen uniformly on compact sets. If
$\mathfrak b\subset\mathfrak a$, then
$\psi_{\mathfrak b}\le\psi_{\mathfrak a}+O(1)$; moreover
$\psi_{\mathfrak a\mathfrak b}=\psi_{\mathfrak a}+\psi_{\mathfrak b}+O(1)$,
and $\psi_{\mathfrak a}\circ\pi=\psi_{\mathfrak a\mathcal O_Y}+O(1)$ for
every modification $\pi\colon Y\to X$. If $Y\subset X$ is a submanifold,
the \emph{restriction} of $\mathfrak a$ to $Y$ is the image ideal
$\rest{Y}\mathfrak a\coloneqq \im(\mathfrak a\to\mathcal O_Y)$
\cite[Definition~1.4.5]{XiaSGPT}; it is generated by the restrictions of
local generators of $\mathfrak a$ (it should not be confused with
$\mathfrak a\otimes\mathcal O_Y$), so that
$\psi_{\rest Y\mathfrak a}=\psi_{\mathfrak a}|_Y+O(1)$.

For an effective divisor $D$ we write
$\psi_D\coloneqq \log|s_D|_{h_D}^2$, where $s_D$ is the canonical section of
$\mathcal O(D)$ and $h_D$ a smooth Hermitian metric on $\mathcal O(D)$;
by the Poincar\'e--Lelong formula, $\ddc\psi_D=[D]-\Theta_D$, where
$\Theta_D$ is the Chern form of $h_D$. For an effective
$\mathbf Q$-divisor we clear denominators. Thus $\psi_D$ is bounded above,
$\psi_{D_1+D_2}=\psi_{D_1}+\psi_{D_2}+O(1)$,
$\psi_{D_1}\le\psi_{D_2}+O(1)$ if $D_1\ge D_2$, and
$\psi_{\mathfrak a}\circ\pi=\psi_D+O(1)$ if
$\mathfrak a\mathcal O_Y=\mathcal O_Y(-D)$ for a modification
$\pi\colon Y\to X$. When exact identities between divisorial potentials or
their Chern forms are used, the metrics are chosen compatibly with sums,
pullbacks and restrictions.

A quasi-psh function $\varphi$ has \emph{analytic singularities of type
$(c,\mathfrak a)$}, where $c\in\mathbf Q_{>0}$ and $\mathfrak a$ is a
coherent ideal, if $\varphi=c\,\psi_\mathfrak a+O(1)$ locally
\cite[Definition~1.6.1]{XiaSGPT}; the remainder is merely bounded.
(Numerical restricted volumes are defined in
\cref{sec:restricted-volumes} with the stronger convention of
\cite[Section~2]{CT22}, with smooth remainder.) If moreover
$\mathfrak a\mathcal O_Y=\mathcal O_Y(-D)$ is invertible on a modification
$\pi\colon Y\to X$ (a \emph{principalization} of $\mathfrak a$), then
$\pi^*\varphi=c\,\psi_D+O(1)$, and we say that $\pi^*\varphi$ has
\emph{log singularities} along the $\mathbf Q$-divisor $cD$. Every
coherent ideal on a compact complex manifold admits a principalization,
and every quasi-psh function with analytic singularities admits a
\emph{log resolution}, a modification on which it has log singularities
along an effective $\mathbf Q$-divisor with simple normal crossing support
\cite[Theorem~1.6.1]{XiaSGPT}. If $\varphi$ has analytic singularities of
type $(c,\mathfrak a)$ and $Y\subset X$ is a submanifold with
$\mathfrak a\not\subset\mathcal I_Y$, then $\varphi|_Y$ has analytic
singularities of type $(c,\rest{Y}\mathfrak a)$.

The cohomological content of log singularities is the following.

\begin{lemma}[{\cite[Lemma~1.6.1]{XiaSGPT}}]\label{lem:residual-nef}
Let $\theta$ be a smooth closed real $(1,1)$-form on a compact K\"ahler
manifold $Z$, and let $\varphi\in\PSH(Z,\theta)$ have log singularities
along an effective $\mathbf Q$-divisor $D$. Then the class $\{\theta\}-\{D\}$ is
nef. If moreover $\theta+\ddc\varphi$ is a K\"ahler current, then
$\{\theta\}-\{D\}$ is K\"ahler.
\end{lemma}

\subsection{Non-pluripolar products, envelopes and comparison of
singularities}\label{sec:npp}

Let $\theta$ be a smooth closed real $(1,1)$-form on $X$ and
$\varphi\in\PSH(X,\theta)$. The non-pluripolar product
$\theta_\varphi^n$, which goes back to Bedford--Taylor \cite{BT87} and
was developed by Guedj--Zeriahi \cite{GZ07} and
Boucksom--Eyssidieux--Guedj--Zeriahi \cite{BEGZ10}, is a positive measure
on $X$ putting no mass on pluripolar sets. More generally, mixed
non-pluripolar products
$\theta_{1,\varphi_1}\wedge\dots\wedge\theta_{n,\varphi_n}$ are well
defined and multilinear \cite[Proposition~2.2.1]{XiaSGPT}. Powers of
positive currents are always understood in the non-pluripolar sense. We
use the following facts.
\begin{enumerate}[label=(\alph*),leftmargin=*,itemsep=.2em,topsep=.3em]
\item\label{it:monotone} (Monotonicity \cite{WN19,DDNL18mono}; see
  \cite[Theorem~2.4.3]{XiaSGPT} for the mixed statement.) If
  $\varphi_j\lessing\psi_j$ for all $j$, then
  $\int_X\theta_{1,\varphi_1}\wedge\dots\wedge\theta_{n,\varphi_n}
  \le\int_X\theta_{1,\psi_1}\wedge\dots\wedge\theta_{n,\psi_n}$.
\item\label{it:invariance} (Invariance \cite[Proposition~2.4.2]{XiaSGPT}.)
  If $\pi\colon Y\to X$ is a modification and $T_1,\dots,T_n$ are
  closed positive $(1,1)$-currents on $X$, then
  $\int_Y\pi^*T_1\wedge\dots\wedge\pi^*T_n=\int_XT_1\wedge\dots\wedge T_n$.
\item\label{it:bigness} (\cite[Proposition~1.22]{BEGZ10}; see
  \cite[Proposition~2.4.1]{XiaSGPT}.) If some $\varphi\in\PSH(X,\theta)$
  satisfies $\int_X\theta_\varphi^n>0$, then $\{\theta\}$ is big.
\end{enumerate}
We write $\PSH(X,\theta)_{>0}$ for the set of $\varphi\in\PSH(X,\theta)$
with $\int_X\theta_\varphi^n>0$, and
$V_\theta\coloneqq \sup\{\varphi\in\PSH(X,\theta):\varphi\le0\}$ for the
potential with minimal singularities.

\subsubsection*{Envelopes}
For $\varphi\in\PSH(X,\theta)$, the \emph{envelope of the singularity
type} of $\varphi$ is
\[
  P_\theta[\varphi]
  \coloneqq \sup{}^*\bigl\{\psi\in\PSH(X,\theta):\psi\le0,\ \psi\lessing\varphi\bigr\}
\]
\cite[Definition~3.1.2]{XiaSGPT}, where $\sup^*$ denotes the upper
semicontinuous regularization; such envelopes were introduced by
Ross--Witt~Nystr\"om \cite{RWN14} and studied in big classes by
Darvas--Di~Nezza--Lu \cite{DDNL18full,DDNL18mono}. For a function
$f\colon X\to[-\infty,\infty]$, the \emph{rooftop envelope} is
\[
  P_\theta(f)\coloneqq \sup{}^*\bigl\{\psi\in\PSH(X,\theta):
     \psi\le f\ \text{quasi-everywhere}\bigr\}
\]
\cite[Definition~2.4.1]{XiaSGPT}, where quasi-everywhere means outside a
pluripolar set. It is either identically $\pm\infty$ or $\theta$-psh, and
$P_\theta(f)\le f$ quasi-everywhere.

\subsubsection*{The \texorpdfstring{$P$}{P}- and
\texorpdfstring{$\mathcal I$}{I}-preorders}
Following \cite[Definition~6.1.1]{XiaSGPT}, we write $\varphi\lessP\psi$
if $P_\theta[\varphi]\le P_\theta[\psi]$ for some (equivalently, every)
$\theta$ with $\varphi,\psi\in\PSH(X,\theta)_{>0}$. We write
$\varphi\lessI\psi$ if
$\mathcal I(\lambda\varphi)\subset\mathcal I(\lambda\psi)$ for all
$\lambda>0$. The corresponding equivalence relations are denoted
$\simP$ and $\simI$, and the preorders apply to closed positive
$(1,1)$-currents through their potentials. We use the following facts.
\begin{enumerate}[label=(\alph*),leftmargin=*,itemsep=.2em,topsep=.3em,resume]
\item\label{it:I-valuative} $\varphi\lessI\psi$ if and only if
  $\nu(\varphi,E)\ge\nu(\psi,E)$ for every prime divisor $E$ over $X$
  \cite[Proposition~6.1.2]{XiaSGPT} (see also \cite{BFJ08}); this
  relation is invariant under modifications
  \cite[Lemma~6.1.4]{XiaSGPT}.
\item\label{it:orders} $\varphi\lessing\psi\Rightarrow\varphi\lessP\psi
  \Rightarrow\varphi\lessI\psi$ (see the remark following
  \cite[Definition~6.1.1]{XiaSGPT}, and
  \cite[Corollary~6.1.1]{XiaSGPT}). In particular, $\mathcal I$-equivalent,
  and a fortiori $P$-equivalent, functions have the same multiplier ideals
  and the same generic Lelong numbers along prime divisors over $X$
  \cite[Proposition~3.2.1]{XiaSGPT}.
\item\label{it:I-envelope} For $\varphi\in\PSH(X,\theta)$, the
  $\mathcal I$-envelope $P_\theta[\varphi]_{\mathcal I}$ of Darvas--Xia
  \cite{DX22} satisfies $\varphi\simI P_\theta[\varphi]_\mathcal I$
  \cite[Proposition~3.2.6]{XiaSGPT}.
\end{enumerate}

The following theorem is the source of the margin described in the
introduction.

\begin{theorem}[{\cite[Theorem~6.1.1]{XiaSGPT}}]\label{thm:rooftop}
Let $\varphi,\psi\in\PSH(X,\theta)_{>0}$ with $\varphi\lessP\psi$, and
let $t>1$ satisfy
$t^n\int_X\theta_\varphi^n>(t^n-1)\int_X\theta_\psi^n$. Then
$P_\theta\bigl(t\varphi-(t-1)\psi\bigr)\in\PSH(X,\theta)_{>0}$.
\end{theorem}

\subsubsection*{The pseudometric \texorpdfstring{$\dS$}{dS}}
For a big class $\{\theta\}$, let $\dS=d_{S,\theta}$ be the pseudometric
on $\PSH(X,\theta)$ of Darvas--Di~Nezza--Lu \cite{DDNL21}; see
\cite[Definition~6.2.1]{XiaSGPT}. We only use the following properties.
\begin{enumerate}[label=(\alph*),leftmargin=*,itemsep=.2em,topsep=.3em,resume]
\item\label{it:dS-zero} $\dS(\varphi,\psi)=0$ if and only if
  $\varphi\simP\psi$ \cite[Proposition~6.2.2]{XiaSGPT}.
\item\label{it:dS-background} For a K\"ahler form $\omega'$, a net
  converges in $d_{S,\theta}$ if and only if it converges in
  $d_{S,\theta+\omega'}$ \cite[Corollary~6.2.9]{XiaSGPT}; thus
  $\dS$-convergence does not depend on the background form.
\item\label{it:dS-mass} If $\theta_1,\dots,\theta_n$ represent big
  classes and $\varphi_j^i\xrightarrow{\dS}\varphi_j$ in
  $\PSH(X,\theta_j)$ for $j=1,\dots,n$, then
  $\int_X\theta_{1,\varphi^i_1}\wedge\dots\wedge\theta_{n,\varphi^i_n}
  \to\int_X\theta_{1,\varphi_1}\wedge\dots\wedge\theta_{n,\varphi_n}$
  \cite[Theorem~6.2.1]{XiaSGPT}.
\item\label{it:dS-closed} If $\varphi_i\xrightarrow{\dS}\varphi$ and
  $\varphi_i\lessP\psi$ for all $i$, then $\varphi\lessP\psi$
  \cite[Proposition~6.2.5]{XiaSGPT}.
\item\label{it:dS-lelong} If $\varphi_i\xrightarrow{\dS}\varphi$, then
  $\nu(\varphi_i,E)\to\nu(\varphi,E)$ for every prime divisor $E$ over
  $X$ \cite[Theorem~6.2.4]{XiaSGPT}, a result originally proved in
  \cite{Xia24}.
\end{enumerate}

\subsubsection*{Quasi-equisingular approximations}
Let $\varphi\in\PSH(X,\theta)$. A sequence $(\varphi_j)_{j\ge1}$ of
quasi-psh functions with analytic singularities is a
\emph{quasi-equisingular approximation} of $\varphi$ if it decreases to
$\varphi$, if $\varphi_j\in\PSH(X,\theta+\epsilon_j\omega)$ for a
sequence $\epsilon_j\downarrow0$, and if for all $\lambda'>\lambda>0$
one has $\mathcal I(\lambda'\varphi_j)\subset\mathcal I(\lambda\varphi)$
for $j$ large \cite[Definition~1.6.4]{XiaSGPT}. By Demailly's
regularization theorem \cite{Dem92,DPS01}, in the form of
\cite[Theorem~1.6.2]{XiaSGPT}, every $\varphi\in\PSH(X,\theta)$ admits a
quasi-equisingular approximation $(\varphi_j)_j$ such that $\varphi_j$ has
analytic singularities of type $(2^{-j},\mathcal I(2^j\varphi))$. If
$\varphi\in\PSH(X,\theta)_{>0}$, a decreasing sequence $(\varphi_j)_j$ of
potentials with analytic singularities, with
$\varphi_j\in\PSH(X,\theta+\epsilon_j\omega)$, $\epsilon_j\downarrow0$ and
$\inf_j\varphi_j=\varphi$, is a quasi-equisingular approximation of
$\varphi$ if and only if $\varphi_j\xrightarrow{\dS}P_\theta[\varphi]_\mathcal I$
\cite[Corollary~7.1.2]{XiaSGPT}; this circle of ideas goes back to
\cite{DX22,DX24b}.

\subsection{Non-K\"ahler locus and non-nef locus}\label{sec:nK}

If $\alpha$ is big and a K\"ahler current $T\in\alpha$ has a potential
which is bounded near a point $x$, then $\nu(T,x)=0$, so
$x\notin\EnK(\alpha)$. By \cite[Theorem~3.17]{Bou04} (see also
\cite[Theorem~1.7.1]{XiaSGPT}), there is a K\"ahler current $R\in\alpha$
with analytic singularities whose polar set is exactly $\EnK(\alpha)$.

For a pseudoeffective class $\alpha$, the \emph{non-nef locus} is
\[
  \Enn(\alpha)\coloneqq \bigcup_{\epsilon>0}\EnK(\alpha+\epsilon\{\omega\});
\]
it does not depend on $\omega$ \cite[Section~2]{CT22}. If $\alpha$ is big
and $\theta\in\alpha$ is a smooth representative, then
$\Enn(\alpha)=\{x\in X:\nu(V_\theta,x)>0\}\subset\EnK(\alpha)$
\cite[Definitions~1.7.6 and~1.7.7]{XiaSGPT}. By Siu's theorem
\cite{Siu74}, for an irreducible analytic subset $Z$ one has
$\nu(V_\theta,Z)=0$ if and only if $Z\not\subset\Enn(\alpha)$. For big
$\alpha$ and a modification $\mu\colon\widetilde X\to X$, we shall use
the transformation rules \cite[Lemmas~4.1 and~4.3]{CT22}
\begin{equation}\label{eq:transformation-rules}
  \Enn(\mu^*\alpha)=\mu^{-1}\bigl(\Enn(\alpha)\bigr),\qquad
  \EnK(\mu^*\alpha)=\mu^{-1}\bigl(\EnK(\alpha)\bigr)\cup\Exc(\mu),
\end{equation}
where $\Exc(\mu)$ is the exceptional locus of $\mu$.

\subsection{Numerical restricted volumes and the trace
operator}\label{sec:restricted-volumes}

Let $\alpha$ be a pseudoeffective class, let $\theta\in\alpha$ be a
smooth representative, and let $V\subset X$ be an irreducible analytic
subset of dimension $k>0$. Following \cite[Section~2]{CT22}, the
\emph{numerical restricted volume} of $\alpha$ along $V$ is
$\langle\alpha^k\rangle_{X\mid V}\coloneqq 0$ if $V\subset\Enn(\alpha)$,
and otherwise
\[
  \langle\alpha^k\rangle_{X\mid V}
  \coloneqq \lim_{\epsilon\downarrow0}\ \sup_{T}
     \int_{V_{\mathrm{reg}}}
     \bigl((T+\epsilon\omega)|_{V_{\mathrm{reg}}}\bigr)_{\mathrm{ac}}^k,
\]
where $T$ runs over the closed real $(1,1)$-currents in $\alpha$ with
$T\ge-\epsilon\omega$ whose polar sets do not contain $V$ and whose local
potentials have the form $c\log\sum_j|f_j|^2+g$ with $c>0$ and $g$ smooth;
the subscript $\mathrm{ac}$ denotes the absolutely continuous part. For
such currents the integrand is the non-pluripolar product, both being the
extension by zero of the smooth product off the polar set
\cite[Remark~2.2]{CT22}. For $V=X$ one recovers the volume,
$\langle\alpha^n\rangle_{X\mid X}=\vol(\alpha)$, which vanishes if and
only if $\alpha$ is not big \cite{Bou02b}. We use two further properties
from \cite{CT22}.
\begin{enumerate}[label=(\alph*),leftmargin=*,itemsep=.2em,topsep=.3em]
\item\label{it:rv-positive} If $V\not\subset\EnK(\alpha)$, then
  $\langle\alpha^k\rangle_{X\mid V}>0$. Indeed, by
  \cite[Lemma~2.7]{CT22} the numerical restricted volume is then the
  supremum of
  $\int_{V_{\mathrm{reg}}}(T|_{V_{\mathrm{reg}}})_{\mathrm{ac}}^k$
  over the K\"ahler currents $T\in\alpha$ with local potentials of the
  form above whose polar sets do not contain $V$, and such currents exist
  when $V\not\subset\EnK(\alpha)$. If $T\ge\delta\omega$, the restricted
  integral is at least $\delta^k\int_V\omega^k>0$.
\item\label{it:rv-birational} If $\alpha$ is big,
  $\mu\colon\widetilde X\to X$ is a modification with
  $V\not\subset\mu(\Exc(\mu))$, and $\widetilde V$ is the strict transform
  of $V$, then $\langle(\mu^*\alpha)^k\rangle_{\widetilde X\mid\widetilde V}
  =\langle\alpha^k\rangle_{X\mid V}$ \cite[Lemma~4.2]{CT22}.
\end{enumerate}

\subsubsection*{The trace operator}
Let $V\subset X$ be a connected submanifold of dimension $k$, and let
$\varphi$ be a quasi-psh function on $X$ with $\nu(\varphi,V)=0$. If
$(\varphi_j)_j$ is a quasi-equisingular approximation of $\varphi$, then
$\varphi_j|_V\not\equiv-\infty$ for all $j$, and the restrictions
$\varphi_j|_V$ converge in $\dS$, with respect to any sufficiently
positive K\"ahler form on $V$, to a quasi-psh function on $V$ which is
independent of the approximation up to $P$-equivalence
\cite[Propositions~8.1.1 and~8.1.2, Definition~8.1.1]{XiaSGPT}. Any such
limit is called a \emph{trace} of $\varphi$ along $V$ and is denoted by
$\trc_V(\varphi)$. If $\varphi\in\PSH(X,\theta)$ and a trace of
$\varphi$ can be chosen in $\PSH(V,\theta|_V)_{>0}$, one sets
\[
  \trc^\theta_V(\varphi)\coloneqq P_{\theta|_V}\bigl[\trc_V(\varphi)\bigr]
  \in\PSH(V,\theta|_V)_{>0}
\]
\cite[Definition~8.1.2]{XiaSGPT}. The trace operator was introduced in
\cite{DX24}; we need two facts about it.

\begin{theorem}[{\cite[Theorem~8.3.1]{XiaSGPT}}]\label{thm:trace-formula}
Let $\alpha$ be a big class, $\theta\in\alpha$ a smooth representative,
and $V\subset X$ a connected submanifold of dimension $k>0$ with
$V\not\subset\Enn(\alpha)$. Then
\[
  \langle\alpha^k\rangle_{X\mid V}
  =\int_V\bigl(\theta|_V+\ddc\trc^\theta_V(V_\theta)\bigr)^k,
\]
where the right-hand side is understood to be $0$ if
$\trc^\theta_V(V_\theta)$ is not defined.
\end{theorem}

\begin{proof}
For $\epsilon>0$, put
$T_\epsilon\coloneqq \theta+\epsilon\omega+\ddc V_{\theta+\epsilon\omega}$.
Since $V\not\subset\Enn(\alpha)\supset\EnK(\alpha+\epsilon\{\omega\})$
and $V_{\theta+\epsilon\omega}$ is locally bounded outside
$\EnK(\alpha+\epsilon\{\omega\})$, the restriction $T_\epsilon|_V$ is well
defined, and by \cite[Remark~2.2 and equation~(2.3)]{CT22},
\[
  \langle\alpha^k\rangle_{X\mid V}
  =\lim_{\epsilon\downarrow0}\int_V(T_\epsilon|_V)^k .
\]
By the equality of the second and the last quantities displayed in
\cite[Theorem~8.3.1]{XiaSGPT}, an identity between non-pluripolar masses
whose proof does not involve the definition of restricted volumes, this
limit equals the right-hand side of the formula, with the same convention
when the trace is not defined.
\end{proof}

\begin{theorem}[{\cite[Proposition~8.2.3]{XiaSGPT}}]\label{thm:trace-multiplier}
Let $V\subset X$ be a connected submanifold and let $\varphi$ be a
quasi-psh function on $X$ with $\nu(\varphi,V)=0$. Then
$\mathcal I\bigl(\lambda\trc_V(\varphi)\bigr)\subset\rest{V}\mathcal I(\lambda\varphi)$
for every $\lambda>0$.
\end{theorem}

Since $P$-equivalent functions have the same multiplier ideals
(\cref{sec:npp}\ref{it:orders}), the inclusion holds for every trace of
$\varphi$, and in particular for $\trc^\theta_V(\varphi)$ when it is
defined.

\section{Auxiliary results}\label{sec:aux}

\subsection{A bigness criterion}

We use the following criterion of Xiao, a weak form of the transcendental
holomorphic Morse inequalities. It is the special case of
\cite[Theorem~1.4]{Xiao15}, stated there for Hermitian metrics $\omega$
with $\partial\bar\partial\omega^j=0$, in which the metric is K\"ahler.
The sharp criterion, with $4k$ replaced by $k$, was proved by Popovici
\cite[Theorem~1.1]{Pop16}, confirming a conjecture of
Boucksom--Demailly--P\u{a}un--Peternell in the K\"ahler case; the weak
form suffices for our purposes.

\begin{theorem}[{\cite[Theorem~1.4]{Xiao15}}]\label{thm:nef-difference}
Let $Z$ be a compact K\"ahler manifold of dimension $k$ and let
$A,B$ be nef classes on $Z$. If $A^k-4kA^{k-1}\cdot B>0$, then $A-B$ is
big.
\end{theorem}

\subsection{Divisorial domination}

The following lemma turns a lower bound on the non-pluripolar mass of a
potential into a K\"ahler current which is at least as singular as the
potential in the sense of multiplier ideals, with a positivity constant
that depends only on the lower bound.

\begin{lemma}\label{lem:domination}
Let $(Z,\omega_Z)$ be a connected compact K\"ahler manifold of
dimension $k$, let $\gamma$ be a smooth closed real $(1,1)$-form on $Z$,
and let $h_0>0$. There is a constant $\delta_*\in(0,1]$, depending only on
$(Z,\omega_Z)$, $\gamma$ and $h_0$, with the following property. If
$0\le\eta\le1$, $\gamma_\eta\coloneqq\gamma+\eta\omega_Z$, and
$v\in\PSH(Z,\gamma_\eta)$ satisfies $\int_Z(\gamma_\eta)_v^k\ge h_0$, then
there is a closed positive current $S\in\{\gamma_\eta\}$ such that
\[
  S\ge\delta_*\omega_Z
  \qquad\text{and}\qquad S\lessI v.
\]
\end{lemma}

\begin{proof}
Choose $C^\dagger>0$ such that $\gamma+C^\dagger\omega_Z$ is K\"ahler, put
\[
  C_*\coloneqq\max_{0\le s\le 2+C^\dagger}
  \bigl(\{\gamma\}+s\{\omega_Z\}\bigr)^{k-1}\cdot\{\omega_Z\},
\]
and choose $\delta_*\in(0,1]$ with $8k\delta_*C_*<h_0$. Let $\eta$ and $v$ be as
in the statement.

\emph{Step 1: analytic approximation.} The class $\{\gamma_\eta\}$ is big
by \cref{sec:npp}\ref{it:bigness}. By \cite[Theorem~1.6.2]{XiaSGPT},
$v$ admits a quasi-equisingular approximation $(v_j)_{j}$ with
$v_j\in\PSH(Z,\gamma_j)$, where $\gamma_j\coloneqq\gamma_\eta+\epsilon_j\omega_Z$
and $\epsilon_j\downarrow0$; we may assume that $\epsilon_j\le\delta_*$ for all $j$.
Since $v\in\PSH(Z,\gamma_\eta)_{>0}$, \cite[Corollary~7.1.2]{XiaSGPT}
gives $v_j\xrightarrow{\dS}P_{\gamma_\eta}[v]_\mathcal I$, hence, by
\cref{sec:npp}\ref{it:dS-lelong}, \ref{it:orders} and~\ref{it:I-envelope},
\begin{equation}\label{eq:dom-lelong-conv}
  \nu(v_j,E)\to\nu(v,E)
  \qquad\text{for every prime divisor $E$ over $Z$}.
\end{equation}
Since $v\le v_j$, monotonicity (\cref{sec:npp}\ref{it:monotone}) and
multilinearity give
$\int_Z(\gamma_j)_{v_j}^k\ge\int_Z\bigl((\gamma_\eta)_v+\epsilon_j\omega_Z\bigr)^k
\ge\int_Z(\gamma_\eta)_v^k\ge h_0$.

\emph{Step 2: a K\"ahler current on a log resolution.} Let
$\rho_j\colon Z_j\to Z$ be a log resolution of $v_j$, and write the Siu
decomposition of the pullback as
\[
  \rho_j^*\bigl((\gamma_j)_{v_j}\bigr)=T_j+[F_j],
\]
where $F_j$ is an effective $\mathbf Q$-divisor and $T_j$ is a closed
positive current with locally bounded potentials. The class
$A_j\coloneqq\{T_j\}=\rho_j^*\{\gamma_j\}-\{F_j\}$ is nef by
\cref{lem:residual-nef}. Since non-pluripolar products do not charge
$F_j$, Bedford--Taylor theory, \cref{sec:npp}\ref{it:invariance} and
Step~1 give
\[
  A_j^k=\int_{Z_j}T_j^k
  =\int_{Z_j}\bigl(\rho_j^*(\gamma_j)_{v_j}\bigr)^k
  =\int_Z(\gamma_j)_{v_j}^k\ge h_0.
\]
The class
$\kappa_j\coloneqq\rho_j^*\{\gamma+(\eta+\epsilon_j+C^\dagger)\omega_Z\}$ is nef
and $\kappa_j-A_j=\{F_j\}+C^\dagger\rho_j^*\{\omega_Z\}$ is pseudoeffective.
Since the product of a pseudoeffective class with nef classes is
nonnegative, expanding
$\kappa_j^{k-1}-A_j^{k-1}=\sum_{i=0}^{k-2}\kappa_j^i\cdot(\kappa_j-A_j)\cdot A_j^{k-2-i}$
gives
\[
  A_j^{k-1}\cdot\rho_j^*\{\omega_Z\}
  \le \kappa_j^{k-1}\cdot\rho_j^*\{\omega_Z\}
  =\bigl(\{\gamma\}+(\eta+\epsilon_j+C^\dagger)\{\omega_Z\}\bigr)^{k-1}
   \cdot\{\omega_Z\}
  \le C_*.
\]
As $\epsilon_j\le\delta_*$, it follows that
$A_j^k-4k(\epsilon_j+\delta_*)A_j^{k-1}\cdot\rho_j^*\{\omega_Z\}
\ge h_0-8k\delta_*C_*>0$. By \cref{thm:nef-difference}, applied on $Z_j$
to the nef classes $A_j$ and $(\epsilon_j+\delta_*)\rho_j^*\{\omega_Z\}$,
the class $A_j-(\epsilon_j+\delta_*)\rho_j^*\{\omega_Z\}$ contains a
K\"ahler current $S'_j$. Put
\[
  \widetilde S_j\coloneqq S'_j+\delta_*\rho_j^*\omega_Z+[F_j].
\]
Then $\widetilde S_j\ge\delta_*\rho_j^*\omega_Z$,
$\widetilde S_j\ge[F_j]$, and
$\{\widetilde S_j\}=A_j+\{F_j\}-\epsilon_j\rho_j^*\{\omega_Z\}
=\rho_j^*\{\gamma_\eta\}$.

\emph{Step 3: descent and passage to the limit.} By the $\ddc$-lemma,
$\widetilde S_j-\delta_*\rho_j^*\omega_Z
=\rho_j^*(\gamma_\eta-\delta_*\omega_Z)+\ddc\widetilde s_j$ with
$\widetilde s_j\in\PSH\bigl(Z_j,\rho_j^*(\gamma_\eta-\delta_*\omega_Z)\bigr)$.
By \cite[Proposition~1.5.3]{XiaSGPT}, $\widetilde s_j=s_j\circ\rho_j$ for
some $s_j\in\PSH(Z,\gamma_\eta-\delta_*\omega_Z)$, and
$S_j\coloneqq\gamma_\eta+\ddc s_j$ satisfies $S_j\ge\delta_*\omega_Z$ and
$\rho_j^*S_j=\widetilde S_j$. Since $T_j$ has locally bounded potentials,
it has vanishing Lelong numbers, so $\widetilde S_j\ge[F_j]$ implies
$\nu(\widetilde S_j,E)\ge\nu([F_j],E)=\nu(\rho_j^*v_j,E)$ for every prime
divisor $E$ over $Z_j$. By \cref{sec:npp}\ref{it:I-valuative}, applied on
$Z_j$ and then on $Z$, we get $\widetilde S_j\lessI\rho_j^*v_j$, then
$S_j\lessI v_j$, and finally $\nu(S_j,E)\ge\nu(v_j,E)$ for every prime
divisor $E$ over $Z$.

Normalize $\sup_Zs_j=0$. By compactness, a subsequence of $(s_j)_j$
converges in $L^1(Z)$ to some $s\in\PSH(Z,\gamma_\eta-\delta_*\omega_Z)$,
and $S\coloneqq\gamma_\eta+\ddc s$ satisfies $S\ge\delta_*\omega_Z$. For
every prime divisor $E$ over $Z$, the upper semicontinuity of Lelong
numbers under weak convergence, applied on a modification on which $E$
lives, and \eqref{eq:dom-lelong-conv} give
\[
  \nu(S,E)\ge\varlimsup_{j\to\infty}\nu(S_j,E)
  \ge\lim_{j\to\infty}\nu(v_j,E)=\nu(v,E).
\]
Hence $S\lessI v$ by \cref{sec:npp}\ref{it:I-valuative}.
\end{proof}

\subsection{Regularization with controlled positivity}

We shall glue local quasi-psh functions by the following elementary
patching procedure, a variant of Richberg's technique \cite{Ric68}. Let
$(Z,\omega_Z)$ be a compact K\"ahler manifold. A \emph{patching datum}
on $Z$ consists of finitely many open sets
$U_j\Subset U_j'\Subset B_j$, $j=1,\dots,J$, such that the $U_j$ cover
$Z$, together with functions $\chi_j\in C^\infty_c(U'_j)$ with
$0\le\chi_j\le1$ and $\chi_j=1$ on $U_j$; we fix $C_\chi>0$ with
$\ddc\chi_j\ge-C_\chi\omega_Z$ for all $j$.

\begin{lemma}\label{lem:patching}
Let $(U_j,U'_j,B_j,\chi_j)_j$ be a patching datum on $Z$, let $\zeta$ be a
smooth closed real $(1,1)$-form on $Z$, and let $a\ge0$. Let $\varphi_j$
be a quasi-psh function on $B_j$ with $\zeta+\ddc\varphi_j\ge0$, for
$j=1,\dots,J$, and assume that $\varphi_j\le\varphi_l+a$ on
$\overline U{}'_j\cap\overline U{}'_l$ for all $j,l$. Then
\[
  \varphi(z)\coloneqq\max\bigl\{\varphi_j(z)+a\chi_j(z):\
  j\text{ such that }z\in\overline U{}'_j\bigr\}
\]
is a quasi-psh function on $Z$ with $\zeta+\ddc\varphi\ge-aC_\chi\omega_Z$,
and $\varphi_j\le\varphi\le\varphi_j+2a$ on $U_j$ for every $j$.
\end{lemma}

\begin{proof}
Fix $z_0\in Z$ and choose $i$ with $z_0\in U_i$. Let $J_0$ be the set of
indices $j$ with $z_0\in U'_j$, and $J_1$ the set of those with
$z_0\in\partial U'_j$. Choose a neighborhood $O$ of $z_0$ contained in
$U_i$ and in $U'_j$ for $j\in J_0$, disjoint from $\overline U{}'_j$ for
$j\notin J_0\cup J_1$, and disjoint from the compact set
$\operatorname{supp}\chi_j\subset U'_j$ for $j\in J_1$. If $z\in O$ and
$j\in J_1$ are such that $z\in\overline U{}'_j$, then $\chi_j(z)=0$ and
$\chi_i(z)=1$, so that
$\varphi_j(z)+a\chi_j(z)\le\varphi_i(z)+a\chi_i(z)$, with $i\in J_0$.
Hence $\varphi=\max_{j\in J_0}(\varphi_j+a\chi_j)$ on $O$. Each function
$\varphi_j+a\chi_j$ with $j\in J_0$ is quasi-psh on $O$ and satisfies
$\zeta+\ddc(\varphi_j+a\chi_j)\ge-aC_\chi\omega_Z$ there, and so does
their maximum. Finally, if $z\in U_j$, the $j$-th function enters the
definition of $\varphi(z)$ with $\chi_j(z)=1$, so $\varphi(z)\ge\varphi_j(z)$;
and every function entering it satisfies
$\varphi_l(z)+a\chi_l(z)\le\varphi_j(z)+2a$, since
$z\in\overline U{}'_l\cap\overline U{}'_j$.
\end{proof}

The next lemma is Demailly's regularization theorem \cite{Dem92,DPS01} in
a quantitative form: the loss of positivity is $O(m^{-1})$, with a constant
that does not depend on the function being regularized. The qualitative
version, with an unspecified loss $\epsilon_m\downarrow0$ but with the
precise singularity type $(m^{-1},\mathcal I(m\varphi))$, is
\cite[Theorem~1.6.2]{XiaSGPT}; it suffices for all the regularizations on
$V$ below. The rate $O(m^{-1})$ is needed exactly once, for the ambient
current of \cref{sec:ambient}. The lemma was stated in \cite{Pop04}, but
we are unable to follow the arguments there, and we include a complete
proof. We only record the lower bound on the singularities of the
regularization, which is all that will be used.

\begin{lemma}\label{lem:regularization}
Let $(Z,\omega_Z)$ be a compact K\"ahler manifold of dimension $d$. There
is a constant $c>0$, depending only on $(Z,\omega_Z)$, with the following
property. For every smooth closed real $(1,1)$-form $\zeta$ on $Z$, every
$\varphi\in\PSH(Z,\zeta)$ and every integer $m\ge1$, there is a quasi-psh
function $\varphi_m$ on $Z$ such that
\begin{enumerate}[label=\textup{(\roman*)},leftmargin=*,itemsep=.2em,topsep=.3em]
\item\label{it:reg-positivity}
  $\zeta+\ddc\varphi_m\ge-\dfrac{c}{m}\,\omega_Z$;
\item\label{it:reg-lower} locally on $Z$,
  $\varphi_m\ge m^{-1}\psi_{\mathcal I(m\varphi)}+O(1)$.
\end{enumerate}
\end{lemma}

\begin{proof}
Fix, once and for all, finitely many coordinate balls covering $Z$. Using a Lebesgue number of this cover, choose coordinate
balls $B_j^+$ with concentric balls $U_j\Subset U'_j\Subset B_j\Subset B_j^+$,
$j=1,\dots,J$, such that the $U_j$ cover $Z$ and such that, whenever
$\overline U{}'_j\cap\overline U{}'_l\ne\varnothing$, the union
$B_j^+\cup B_l^+$ is contained in one of these balls, which we
denote by $\mathcal B_{jl}$. Complete these sets to a patching datum by
choosing functions $\chi_j$. On the finitely many balls involved, the
volume form $\mathrm{d}V\coloneqq\omega_Z^d/d!$ is uniformly comparable
with the Lebesgue measures of the coordinates. All constants below depend
only on these choices.

\emph{Step 1: local Bergman kernels.} On each $B_j^+$ choose a smooth
function $h_j$ with $\ddc h_j=\zeta$; then $\varphi+h_j$ is psh on
$B_j^+$. For $m\ge1$, let $\mathcal H_{j}$ be the Hilbert space of
holomorphic functions $f$ on $B_j$ with
$\|f\|^2\coloneqq\int_{B_j}|f|^2\mathrm{e}^{-m(\varphi+h_j)}\,\mathrm{d}V<\infty$,
let $\mathcal K_j\coloneqq\sup\{|f|^2:f\in\mathcal H_j,\ \|f\|\le1\}$ be its
Bergman kernel on the diagonal, which is the sum of the squared moduli of
the elements of an orthonormal basis, and set
\[
  \varphi_{j,m}\coloneqq\frac1m\log\mathcal K_j-h_j .
\]
Then $\zeta+\ddc\varphi_{j,m}=m^{-1}\ddc\log\mathcal K_j\ge0$ on $B_j$. Since
$\varphi+h_j$ is psh on $B_j^+\Supset B_j$, the ideal
$\mathcal I(m\varphi)=\mathcal I(m(\varphi+h_j))$ is generated on $B_j$ by
elements of $\mathcal H_j$ \cite[Proposition~1.2]{DEL00} (applied with
the weight $\tfrac m2(\varphi+h_j)$, the multiplier ideal being defined
there with the factor $\mathrm{e}^{-2\varphi}$). By coherence, near a
given point of $B_j$ the ideal $\mathcal I(m\varphi)$ is generated by
finitely many $f_1,\dots,f_s\in\mathcal H_j$ of norm at most $1$, and
$\mathcal K_j\ge s^{-1}\sum_i|f_i|^2$ there. In particular $\varphi_{j,m}$ is not
identically $-\infty$, and
\begin{equation}\label{eq:reg-local-lower}
  \varphi_{j,m}\ge m^{-1}\psi_{\mathcal I(m\varphi)}+O(1)
  \qquad\text{locally on }B_j .
\end{equation}

\emph{Step 2: comparison on overlaps.} We claim that there is a constant
$C_B$, independent of $\zeta$, $\varphi$ and $m$, such that
\begin{equation}\label{eq:reg-overlap}
  \varphi_{j,m}\le\varphi_{l,m}+\frac{C_B}m
  \qquad\text{on }\overline U{}'_j\cap\overline U{}'_l .
\end{equation}
Fix $j,l$ with $\overline U{}'_j\cap\overline U{}'_l\ne\varnothing$, and
work in the coordinates $z$ of $\mathcal B_{jl}$. Choose a smooth function $g$
on $\mathcal B_{jl}$ with $\ddc g=\zeta$. The functions $h_j-g$ and $h_l-g$ are
pluriharmonic on the balls $B_j^+$ and $B_l^+$, so $h_j-g=\Re a_j$ and
$h_l-g=\Re a_l$ with $a_j$, $a_l$ holomorphic. Put $\Psi\coloneqq\varphi+g$,
a psh function on $\mathcal B_{jl}$, and for $B\in\{B_j,B_l\}$ let $\mathcal K^\Psi_B$
be the Bergman kernel on the diagonal of the space of holomorphic
functions on $B$ which are square integrable against
$\mathrm{e}^{-m\Psi}\mathrm{d}V$. Since
$\mathrm{e}^{-m(\varphi+h_j)}=\bigl|\mathrm{e}^{-ma_j/2}\bigr|^2\mathrm{e}^{-m\Psi}$,
multiplication by $\mathrm{e}^{-ma_j/2}$ is an isometry from $\mathcal H_j$
onto this space for $B=B_j$, so that
$\mathcal K^\Psi_{B_j}=\mathcal K_j\,\mathrm{e}^{-m\Re a_j}$, that is,
$\varphi_{j,m}=m^{-1}\log\mathcal K^\Psi_{B_j}-g$; similarly
$\varphi_{l,m}=m^{-1}\log\mathcal K^\Psi_{B_l}-g$. It therefore suffices to find
$C_0$, independent of $\Psi$, $m$ and $x$, such that for every
$x\in\overline U{}'_j\cap\overline U{}'_l$ and every $f$ in the unit ball
of the $\Psi$-weighted space on $B_j$ there is a holomorphic function $F$
on $B_l$ with $F(x)=f(x)$ and $\Psi$-weighted norm at most $C_0$; then
$\mathcal K^\Psi_{B_l}(x)\ge C_0^{-2}\mathcal K^\Psi_{B_j}(x)$, and \eqref{eq:reg-overlap}
holds with $C_B=2\log C_0$ for this pair, and in general with the maximum
over the finitely many pairs.

Choose $\varrho\in C^\infty_c(B_j\cap B_l)$ with $0\le\varrho\le1$ and
$\varrho=1$ near $\overline U{}'_j\cap\overline U{}'_l$, let $r>0$ be the
distance from $\overline U{}'_j\cap\overline U{}'_l$ to the support of
$\bar\partial\varrho$. Fix $x$
and $f$ as above, extend $\varrho f$ by zero to $B_l$, and solve
$\bar\partial u=\bar\partial(\varrho f)$ on the pseudoconvex domain $B_l$
with respect to the weight $m\Psi+2d\log|z-x|+|z|^2$, whose complex
Hessian dominates that of $|z|^2$. H\"ormander's $L^2$-estimate
\cite[Chapter~VIII, \S6]{DemCADG} gives a solution with
\[
  \int_{B_l}|u|^2\frac{\mathrm{e}^{-m\Psi-|z|^2}}{|z-x|^{2d}}\,\mathrm{d}V
  \le C_H\int_{B_l}|\bar\partial\varrho|^2|f|^2
  \frac{\mathrm{e}^{-m\Psi-|z|^2}}{|z-x|^{2d}}\,\mathrm{d}V
  \le C_Hr^{-2d}\sup|\bar\partial\varrho|^2 ,
\]
where $C_H$ depends only on $d$ and on the comparison between
$\mathrm{d}V$ and the Lebesgue measure. Since $|z-x|$ and $|z|$ are bounded
on $B_l$ by the diameter of $\mathcal B_{jl}$ and a bound for the
coordinates, the $\Psi$-weighted norm of $u$ is bounded by a constant $C_1$. Moreover $u$
is holomorphic near $x$, where $\bar\partial\varrho=0$, and $\Psi$ is
bounded above near $x$; the local integrability of $|u|^2|z-x|^{-2d}$
therefore forces $u(x)=0$. Hence $F\coloneqq\varrho f-u$ is holomorphic on
$B_l$, $F(x)=f(x)$, and $F$ has $\Psi$-weighted norm at most $C_0\coloneqq
1+C_1$. This proves \eqref{eq:reg-overlap}.

\emph{Step 3: patching.} By \eqref{eq:reg-overlap}, \cref{lem:patching}
applies to the functions $\varphi_{j,m}$ with $a=C_B/m$. The resulting
function $\varphi_m$ satisfies
$\zeta+\ddc\varphi_m\ge-(C_BC_\chi/m)\omega_Z$, which is
\ref{it:reg-positivity} with $c\coloneqq C_BC_\chi$, and
$\varphi_m\ge\varphi_{j,m}$ on $U_j$, which gives \ref{it:reg-lower} by
\eqref{eq:reg-local-lower}.
\end{proof}

The same patching produces global quasi-psh functions with prescribed
singularities.

\begin{lemma}\label{lem:ideal-function}
Let $Z$ be a connected compact K\"ahler manifold and let
$\mathfrak a\subset\mathcal O_Z$ be a nonzero coherent ideal. There is a
quasi-psh function $\Psi_{\mathfrak a}$ on $Z$ with
$\Psi_{\mathfrak a}=\psi_{\mathfrak a}+O(1)$ locally. In particular,
$\Psi_{\mathfrak a}$ is locally bounded outside the zero locus of
$\mathfrak a$.
\end{lemma}

\begin{proof}
Choose a patching datum $(U_j,U'_j,B_j,\chi_j)_j$ such that $\mathfrak a$
is generated on a neighborhood of $\overline B_j$ by finitely many
sections $g_{j1},\dots,g_{jN_j}$, and put
$\psi_j\coloneqq\log\sum_i|g_{ji}|^2$, a psh function on $B_j$ which is
not identically $-\infty$ because $\mathfrak a\ne0$ and $Z$ is connected.
Since both families generate $\mathfrak a$ near the compact set
$\overline U{}'_j\cap\overline U{}'_l$, there is a constant $C_B$ with
$\psi_j\le\psi_l+C_B$ there, for all $j,l$. \Cref{lem:patching}, with
$\zeta=0$ and $a=C_B$, yields a quasi-psh function $\Psi_\mathfrak a$ with
$\psi_j\le\Psi_\mathfrak a\le\psi_j+2C_B$ on $U_j$.
\end{proof}

\subsection{Relative principalization}

Let $Y$ be a complex manifold and $W\subset Y$ a connected submanifold. A
composition $f\colon Y'\to Y$ of blowups with smooth centers is
\emph{adapted to $W$} if each center is either a proper analytic subset of
the strict transform of $W$ at that stage, or disjoint from it. By
induction on the number of blowups, the image in $Y$ of each center meets
$W$ in a proper analytic subset. Consequently, $f$ is an isomorphism over
the complement of an analytic subset of $Y$ which meets $W$ in a proper
analytic subset; the strict transform $W'$ of $W$ is a connected
submanifold of $Y'$ mapping onto $W$; and no $f$-exceptional prime
divisor contains $W'$. If $g\colon Y''\to Y'$ is adapted to $W'$, then
$f\circ g$ is adapted to $W$.

\begin{lemma}\label{lem:relative-principalization}
Let $Y$ be a compact complex manifold, let $W\subsetneq Y$ be a
connected submanifold, and let $\mathfrak a\subset\mathcal O_Y$ be a
coherent ideal with $\rest W\mathfrak a\ne0$, that is,
$\mathfrak a\not\subset\mathcal I_W$. There are a composition of blowups
$f\colon Y'\to Y$ adapted to $W$, with strict transform $W'$ of $W$, and an
effective divisor $H$ on $Y'$ supported on the exceptional locus of $f$,
such that
\begin{enumerate}[label=\textup{(\roman*)},leftmargin=*,itemsep=.2em,topsep=.3em]
\item\label{it:rp-ideal-W}
  $\mathcal I_W\mathcal O_{Y'}=\mathcal O_{Y'}(-H)\cdot\mathcal I_{W'}$;
\item\label{it:rp-principal}
  $\mathfrak a\mathcal O_{Y'}=\mathcal O_{Y'}(-H-H')$ for an effective
  divisor $H'$ on $Y'$ disjoint from $W'$;
\item\label{it:rp-restriction}
  $(\rest{W}\mathfrak a)\mathcal O_{W'}=\mathcal O_{W'}(-H|_{W'})$.
\end{enumerate}
\end{lemma}

In words: after the blowups, $\mathfrak a$ is invertible, and near $W'$ it
coincides with the ideal $\mathcal O_{Y'}(-H)$ which
\ref{it:rp-ideal-W} attaches to the blowups alone. Since no component of
$H$ contains $W'$, property~\ref{it:rp-ideal-W} determines $H$. If
$g\colon Y''\to Y'$ is a further composition of blowups adapted to $W'$,
with divisor $H''$ satisfying \ref{it:rp-ideal-W} for $W'$, then
$g^*H+H''$ satisfies \ref{it:rp-ideal-W} for $f\circ g$ and $W$.

In the proof, we use the terminology of \cite{BM08,Wlo05}.

\begin{proof}
Put $\mathfrak b\coloneqq \rest W\mathfrak a\ne0$. By \cite[Theorem~1.3]{BM08}
there is a resolution of the marked ideal $(\mathfrak b,1)$ on $W$: a
finite sequence of blowups of $W$ with smooth centers, each contained
in the cosupport of the current controlled transform of $\mathfrak b$,
after which the controlled transform is the unit ideal. Blowing up the
same centers in the ambient manifold produces a sequence
$f_i\colon Y_i\to Y_{i-1}$, $i=1,\dots,t$, with $Y_0=Y$, whose
$i$-th center $C_i$ is a proper analytic subset of the strict
transform $W_{i-1}$ of $W$. (If $C_i$ has codimension one in
$W_{i-1}$, then $W_i\to W_{i-1}$ is an isomorphism and $E_i|_{W_i}=C_i$.)
Let $E_i\subset Y_i$ be the exceptional divisor, and define inductively
\[
  H_0=0,\quad H_i=f_i^*H_{i-1}+E_i,\qquad
  \mathfrak a_0=\mathfrak a,\quad
  \mathfrak a_i=\mathcal O_{Y_i}(E_i)\,\mathfrak a_{i-1}\mathcal O_{Y_i}.
\]
We show by induction on $i$ that $\mathfrak a_i\subset\mathcal O_{Y_i}$ is
an ideal and that
\begin{equation}\label{eq:rp-induction}
  \mathcal I_W\mathcal O_{Y_i}=\mathcal O_{Y_i}(-H_i)\,\mathcal I_{W_i},\qquad
  \mathfrak a\mathcal O_{Y_i}=\mathcal O_{Y_i}(-H_i)\,\mathfrak a_i,\qquad
  \rest{W_i}\mathfrak a_i=\mathfrak b_i,
\end{equation}
where $\mathfrak b_i$ denotes the controlled transform of $\mathfrak b$.
This is clear for $i=0$. Assume it for $i-1$. By the third identity, the
restriction to $W_{i-1}$ of a local section of $\mathfrak a_{i-1}$ lies in
$\mathfrak b_{i-1}$, hence vanishes on $C_i$, which lies in the
cosupport of $\mathfrak b_{i-1}$. Thus $\mathfrak a_{i-1}\subset\mathcal
I_{C_i}$, so $\mathfrak a_{i-1}\mathcal O_{Y_i}\subset\mathcal
O_{Y_i}(-E_i)$ and $\mathfrak a_i$ is an ideal. The first identity for $i$
follows from
$\mathcal I_{W_{i-1}}\mathcal O_{Y_i}=\mathcal O_{Y_i}(-E_i)\,\mathcal I_{W_i}$,
a computation in coordinates adapted to $C_i\subset W_{i-1}\subset Y_{i-1}$;
the second follows by pulling back its predecessor and inserting the
definition of $\mathfrak a_i$; and restricting the definition of
$\mathfrak a_i$ to $W_i$ gives
$\rest{W_i}\mathfrak a_i=\mathcal O_{W_i}(E_i|_{W_i})\,\mathfrak b_{i-1}\mathcal O_{W_i}
=\mathfrak b_i$.

At the end, $\mathfrak b_t=\mathcal O_{W_t}$, so $\mathfrak a_t$ is the unit
ideal on a neighborhood of $W_t$. Principalize $\mathfrak a_t$ by a
second application of \cite[Theorem~1.3]{BM08}; all its centers lie in
the cosupport of $\mathfrak a_t$, hence are disjoint from the strict
transforms of $W_t$. If $H$ denotes the total transform of $H_t$ and
$H'$ the divisor with $\mathfrak a_t\mathcal O_{Y'}=\mathcal O_{Y'}(-H')$,
then \eqref{eq:rp-induction} gives \ref{it:rp-ideal-W} and
\ref{it:rp-principal}, with $H'\cap W'=\varnothing$. Restricting
\ref{it:rp-principal} to $W'$ gives \ref{it:rp-restriction}. All centers
are adapted to $W$.
\end{proof}

\subsection{K\"ahler forms on blowups}

\begin{lemma}\label{lem:exceptional-metric}
Let $(Z,\omega_Z)$ be a compact K\"ahler manifold and let
$\rho\colon Z'\to Z$ be a composition of blowups with smooth centers.
There are an effective $\rho$-exceptional $\mathbf Q$-divisor $E_\rho$
on $Z'$ and a smooth Hermitian metric on $\mathcal O_{Z'}(E_\rho)$ such that
$\rho^*\omega_Z-\Theta_{E_\rho}$ is a K\"ahler form on $Z'$.
\end{lemma}

\begin{proof}
Blowups along centers of codimension one are isomorphisms and may be
omitted. Consider a single blowup $f\colon Z_1\to Z$ with smooth center
$C$ of codimension at least two and exceptional divisor $F$. By
\cite[Chapter~VII, Proposition~12.4 and its proof]{DemCADG}, $F=\mathbf
P(N_{C/Z})$ and $\mathcal O_{Z_1}(-F)|_F=\mathcal O_{\mathbf P(N_{C/Z})}(1)$,
so a Hermitian metric on $N_{C/Z}$ induces a metric on
$\mathcal O_{Z_1}(-F)|_F$ whose curvature is positive along the fibers of
$F\to C$; extend it smoothly to $\mathcal O_{Z_1}(-F)$. The form
$-\Theta_F$ is then positive on the nonzero vectors of $\ker\mathrm df$,
which are the vectors tangent to the fibers of $F\to C$, while
$f^*\omega_Z$ is semipositive with kernel $\ker\mathrm df$. By compactness
of the unit sphere bundle, $f^*\omega_Z-\epsilon\Theta_F$ is a K\"ahler
form for every sufficiently small $\epsilon>0$. In general, argue by
induction: if $\omega_{i-1}=\rho_{i-1}^*\omega_Z-\Theta_{E_{i-1}}$ is a
K\"ahler form on $Z_{i-1}$ and $f_i\colon Z_i\to Z_{i-1}$ is the next
blowup, with exceptional divisor $F_i$, then
$f_i^*\omega_{i-1}-\epsilon_i\Theta_{F_i}$ is K\"ahler for a small rational
$\epsilon_i>0$, and it equals $\rho_i^*\omega_Z-\Theta_{E_i}$ with
$E_i\coloneqq f_i^*E_{i-1}+\epsilon_iF_i$ and the induced metric.
\end{proof}

\subsection{Gluing}

The following proposition isolates the maximum construction behind the
collar condition. It is used twice: for \cref{lem:no-isolated-points}
below, and at the end of the proof of \cref{thm:main}.

\begin{proposition}\label{prop:gluing}
Let $Y$ be a compact K\"ahler manifold with a smooth closed real
$(1,1)$-form $\theta_Y$, let $W\subset Y$ be a compact subset with an
open neighborhood $U$, and let $y\in W$. Let $\varphi_G$ be a quasi-psh
function on $Y$ and $\varphi_U$ a quasi-psh function on $U$, and assume
that, for some K\"ahler form $\omega_0$ on $Y$,
\begin{enumerate}[label=\textup{(\roman*)},leftmargin=*,itemsep=.25em]
\item\label{it:gl-current}
  $\theta_Y+\ddc\varphi_G\ge\omega_0$ on $Y$, and $\varphi_G(y)=-\infty$;
\item\label{it:gl-local}
  $\theta_Y+\ddc\varphi_U\ge\omega_0$ on $U$, and $\varphi_U$ is bounded
  near $y$;
\item\label{it:gl-collar} for every compact set $\Lambda\subset U\setminus W$
  there is a constant $C_\Lambda$ with $\varphi_U\le\varphi_G+C_\Lambda$
  on $\Lambda$.
\end{enumerate}
Then $\{\theta_Y\}$ contains a K\"ahler current whose potential is bounded
near $y$. In particular, $y\notin\EnK(\{\theta_Y\})$.
\end{proposition}

\begin{proof}
Choose open sets $W\subset U_1\Subset U_2\Subset U$, and let
$c_1\coloneqq C_\Lambda$ for $\Lambda\coloneqq\overline U_2\setminus U_1$. The
function
\[
  \Phi\coloneqq
  \begin{cases}
    \max\{\varphi_U-c_1,\varphi_G\}&\text{on }U_2,\\
    \varphi_G&\text{on }Y\setminus\overline U_1
  \end{cases}
\]
is well defined, because $\varphi_U-c_1\le\varphi_G$ on
$U_2\setminus\overline U_1$ by \ref{it:gl-collar}. By
\ref{it:gl-current} and \ref{it:gl-local}, $\Phi$ is locally the maximum
of $(\theta_Y-\omega_0)$-psh functions, so
$\theta_Y+\ddc\Phi\ge\omega_0$ is a K\"ahler current in $\{\theta_Y\}$.
Since $\varphi_U$ is bounded near $y$, while $\varphi_G(y)=-\infty$ and
$\varphi_G$ is upper semicontinuous, $\Phi=\varphi_U-c_1$ near $y$. The last
assertion follows from \cref{sec:nK}.
\end{proof}

As a first application, we record the following fact, extracted from the
proof of \cite[Lemma~3.1]{CT15}.

\begin{lemma}\label{lem:no-isolated-points}
The non-K\"ahler locus of a big class on a compact K\"ahler manifold has
no isolated points.
\end{lemma}

\begin{proof}
Let $\alpha$ be a big class on a compact K\"ahler manifold $Y$, with a
smooth representative $\theta_Y$, and suppose that $x$ is an isolated
point of $\EnK(\alpha)$. By \cref{sec:nK}, there is a K\"ahler current
$\theta_Y+\ddc\varphi_G$ with analytic singularities whose polar set is
$\EnK(\alpha)$; thus $\varphi_G(x)=-\infty$, $\varphi_G$ is locally
bounded on $Y\setminus\EnK(\alpha)$, and $\theta_Y+\ddc\varphi_G\ge\omega_0$
for some K\"ahler form $\omega_0$. Choose a coordinate ball $B$ centered at
$x$, relatively compact in a coordinate chart, whose closure meets
$\EnK(\alpha)$ only at $x$, and $A>0$ so large that
$\theta_Y+\ddc(A|z|^2)\ge\omega_0$ on $B$. Then
\cref{prop:gluing} applies with $W=\{x\}=\{y\}$, $U=B$ and
$\varphi_U=A|z|^2$: hypothesis~\ref{it:gl-collar} holds because
$\varphi_U$ is bounded above on $B$ and $\varphi_G$ is bounded below on
every compact subset of $B\setminus\{x\}$. Hence
$x\notin\EnK(\alpha)$, a contradiction.
\end{proof}

\section{Reduction to a statement about the trace}\label{sec:setup}

\subsection{Reductions}

Let $X$, $\alpha$ and $V$ be as in \cref{thm:main}, and put
$k\coloneqq \dim V>0$. We have to prove that
$\langle\alpha^k\rangle_{X\mid V}=0$. If $\alpha$ is not big, then
$\EnK(\alpha)=X$, so $V=X$ and
$\langle\alpha^n\rangle_{X\mid X}=\vol(\alpha)=0$
(\cref{sec:restricted-volumes}). We may therefore assume that $\alpha$ is
big; then $\EnK(\alpha)$ is a proper analytic subset of $X$, so $k<n$. If
$V\subset\Enn(\alpha)$, the vanishing holds by definition. We may
therefore assume that $V\not\subset\Enn(\alpha)$, and we argue by
contradiction, assuming that $h\coloneqq \langle\alpha^k\rangle_{X\mid V}>0$.

\begin{lemma}\label{lem:resolution}
In this situation, there is a modification
$\mu\colon\widetilde X\to X$ which is an isomorphism over
$X\setminus V_{\mathrm{sing}}$ and such that the strict transform
$\widetilde V$ of $V$ is smooth. Moreover, $\widetilde X$ is a connected
compact K\"ahler manifold, $\mu^*\alpha$ is big, $\widetilde V$ is an
irreducible component of $\EnK(\mu^*\alpha)$ not contained in
$\Enn(\mu^*\alpha)$, and
$\langle(\mu^*\alpha)^k\rangle_{\widetilde X\mid\widetilde V}=h$.
\end{lemma}

\begin{proof}
Let $\mu$ be an embedded resolution of $V\subset X$, that is, a
composition of blowups with smooth centers lying over $V_{\mathrm{sing}}$,
which is an isomorphism over $X\setminus V_{\mathrm{sing}}$ and such that
the strict transform of $V$ is smooth; it exists by Hironaka's theorem,
see for instance \cite{BM97}. Blowups of compact K\"ahler manifolds along
smooth centers are compact K\"ahler, and the pullback of a big class under
a modification is big. Since $V\not\subset\mu(\Exc(\mu))\subset
V_{\mathrm{sing}}$, \cref{sec:restricted-volumes}\ref{it:rv-birational}
gives $\langle(\mu^*\alpha)^k\rangle_{\widetilde X\mid\widetilde V}=h$. By
\eqref{eq:transformation-rules}, $\Enn(\mu^*\alpha)=\mu^{-1}(\Enn(\alpha))$,
so $\widetilde V\subset\Enn(\mu^*\alpha)$ would give
$V=\mu(\widetilde V)\subset\Enn(\alpha)$, which is excluded.

It remains to see that $\widetilde V$ is an irreducible component of
$\EnK(\mu^*\alpha)=\mu^{-1}(\EnK(\alpha))\cup\Exc(\mu)$. Clearly
$\widetilde V\subset\EnK(\mu^*\alpha)$. Let $Z\supset\widetilde V$ be an
irreducible component of $\EnK(\mu^*\alpha)$. Then $\mu(Z)$ is an
irreducible analytic subset of
$\EnK(\alpha)\cup V_{\mathrm{sing}}=\EnK(\alpha)$ containing $V$, so
$\mu(Z)=V$ because $V$ is an irreducible component of $\EnK(\alpha)$.
Hence $Z\subset\mu^{-1}(V)=\widetilde V\cup\mu^{-1}(V_{\mathrm{sing}})$,
and since $Z$ is irreducible and $\mu(Z)\not\subset V_{\mathrm{sing}}$, we
conclude that $Z=\widetilde V$.
\end{proof}

Replacing $(X,\alpha,V)$ by $(\widetilde X,\mu^*\alpha,\widetilde V)$,
we are reduced to the following situation, which remains in force until
the end of the proof.

\begin{setup}\label{setup}
$X$ is a connected compact K\"ahler manifold of dimension $n$ with
K\"ahler form $\omega$; $\alpha$ is a big class with a smooth
representative $\theta$; $V\subset X$ is a connected submanifold of
dimension $k$, $0<k<n$, which is an irreducible component of
$\EnK(\alpha)$ and is not contained in $\Enn(\alpha)$; and
$h=\langle\alpha^k\rangle_{X\mid V}>0$. We put $\omega_V\coloneqq \omega|_V$
and $\gamma\coloneqq \theta|_V$.
\end{setup}

We shall reach a contradiction in \cref{sec:conclusion}.

\subsection{The trace and its approximations}\label{sec:trace-approx}

Since $V\not\subset\Enn(\alpha)$, we have $\nu(V_\theta,V)=0$, and there
is a point $x_*\in V$ with $\nu(V_\theta,x_*)=0$ (\cref{sec:nK}). By
\cref{thm:trace-formula},
\[
  h=\int_V\bigl(\gamma+\ddc\trc^\theta_V(V_\theta)\bigr)^k,
\]
where the right-hand side would vanish if $\trc^\theta_V(V_\theta)$ were
not defined. Since $h>0$, the trace is defined, and we put
\[
  u\coloneqq \trc^\theta_V(V_\theta)\in\PSH(V,\gamma)_{>0},
  \qquad\int_V\gamma_u^k=h .
\]
For every integer $q\ge1$ we consider the ambient multiplier ideals and
their restrictions to $V$,
\[
  \mathfrak B_q\coloneqq \mathcal I(qV_\theta)\subset\mathcal O_X,
  \qquad
  \mathfrak b_q\coloneqq \rest{V}\mathfrak B_q\subset\mathcal O_V .
\]
By Skoda's theorem and $\nu(V_\theta,x_*)=0$, $\mathfrak B_q$ is the unit
ideal near $x_*$; in particular $\mathfrak b_q\ne0$. By
\cref{thm:trace-multiplier} and the subadditivity
\eqref{eq:subadditivity}, for all integers $q,r\ge1$,
\begin{equation}\label{eq:trace-ideals}
  \mathcal I(qu)\subset\mathfrak b_q,
  \qquad
  \mathfrak B_{rq}\subset\mathfrak B_q^{\,r},
  \qquad
  \mathfrak b_{rq}\subset\mathfrak b_q^{\,r}.
\end{equation}

Let $(\Phi_q)_q$ be a quasi-equisingular approximation of $V_\theta$ as
in \cref{sec:npp}, indexed by the dyadic integers $q=2^j$. Thus $\Phi_q$
decreases to $V_\theta$ as $q\to\infty$, and
\[
  \Phi_q\in\PSH(X,\theta+\eta_q\omega),\qquad
  \eta_q\downarrow0,\qquad
  \Phi_q\ \text{has analytic singularities of type }(q^{-1},\mathfrak B_q).
\]
Since $\mathfrak B_q\not\subset\mathcal I_V$, the restriction
$\phi_q\coloneqq \Phi_q|_V$ is not identically $-\infty$; it belongs to
$\PSH(V,\beta_q)$, where $\beta_q\coloneqq \gamma+\eta_q\omega_V$, and
\begin{equation}\label{eq:phiq-singularities}
  \phi_q=q^{-1}\psi_{\mathfrak b_q}+O(1)\qquad\text{locally on }V.
\end{equation}
By the definition of the trace and \cref{sec:npp}\ref{it:dS-zero},
$\phi_q\xrightarrow{\dS}u$.

\begin{lemma}\label{lem:trace-approximants}
As $q\to\infty$ along the dyadic integers,
\begin{enumerate}[label=\textup{(\roman*)},leftmargin=*,itemsep=.2em,topsep=.3em]
\item\label{it:vq-mass} $\int_V(\beta_q)_{\phi_q}^k\to h$;
\item\label{it:vq-order} $u\lessP\phi_q$ for every sufficiently large $q$.
\end{enumerate}
\end{lemma}

\begin{proof}
Choose $A>\sup_q\eta_q$ such that $\beta\coloneqq \gamma+A\omega_V$ is a
K\"ahler form on $V$, so that $\phi_q\in\PSH(V,\beta)$ for all $q$. Writing
$(\beta_q)_{\phi_q}=\beta_{\phi_q}-(A-\eta_q)\omega_V$ and expanding by
multilinearity, \ref{it:vq-mass} follows from the continuity of mixed
masses (\cref{sec:npp}\ref{it:dS-mass}) and $\eta_q\to0$.

By \ref{it:vq-mass}, $\int_V\beta_{\phi_q}^k\ge\int_V(\beta_q)_{\phi_q}^k>0$
for all large $q$, and $\int_V\beta_u^k\ge\int_V\gamma_u^k=h>0$; so $u$ and
$\phi_q$ lie in $\PSH(V,\beta)_{>0}$ for large $q$. Fix such a $q$. For
dyadic $q'\ge q$ we have $\phi_{q'}\le\phi_q$, hence
$\phi_{q'}\lessP\phi_q$ (\cref{sec:npp}\ref{it:orders}). Since
$\phi_{q'}\xrightarrow{\dS}u$, \cref{sec:npp}\ref{it:dS-closed} gives
$u\lessP\phi_q$.
\end{proof}

\section{A K\"ahler current on the submanifold}\label{sec:positive-current}

We keep \cref{setup} and the notation of \cref{sec:setup}. In this section
we construct a K\"ahler current on $V$ whose multiplier ideals are
controlled by the restricted ideals $\mathfrak b_m$ and $\mathfrak b_M$ at
two levels $m\ll M$.

\subsection{Data fixed first}\label{sec:fixed-data}

Let $\delta_*\in(0,1]$ be the constant of \cref{lem:domination} for
$(Z,\omega_Z,\gamma,h_0)=(V,\omega_V,\gamma,2^{-k-1}h)$, and put
$\delta\coloneqq \delta_*/8$.

By \cref{sec:nK}, there is a K\"ahler current with analytic singularities
\[
  R=\theta+\ddc\Phi_R\in\alpha,\qquad R\ge\delta_R\,\omega,
\]
whose polar set is $\EnK(\alpha)$, for some $\delta_R>0$. Let $Z'$ be
the union of the irreducible components of $\EnK(\alpha)$ other than $V$,
and let $\mathfrak j\coloneqq \mathcal I_{Z'}$, with $\mathfrak j=\mathcal
O_X$ if $Z'=\varnothing$. Since $V\not\subset Z'$, we have
$\rest V\mathfrak j\ne0$. Write $\Phi_R=c_R\,\psi_{\mathfrak r}+O(1)$ locally.
The zero locus of $\mathfrak r$ is $\EnK(\alpha)=V\cup Z'$, so by
R\"uckert's Nullstellensatz \cite[Chapter~II, Theorem~4.22]{DemCADG} and
compactness, $\mathcal I_V^N\mathfrak j^N\subset\mathfrak r$ for some
integer $N$. With $A_R\coloneqq c_RN$, this gives
\begin{equation}\label{eq:PhiR-lower}
  \Phi_R\ge A_R\bigl(\psi_{\mathcal I_V}+\psi_{\mathfrak j}\bigr)+O(1)
  \qquad\text{locally on }X.
\end{equation}

Let $c_0$ be the constant of \cref{lem:regularization} for
$(Z,\omega_Z)=(X,\omega)$, and choose integers $L,C\ge1$ with
\begin{equation}\label{eq:constants}
  L\delta_R>c_0,\qquad C\ge 1+LA_R .
\end{equation}

\subsection{A rooftop envelope with positive mass}

\begin{lemma}\label{lem:positive-rooftop}
For every sufficiently large dyadic integer $m$, there is a function
$v\in\PSH(V,\beta_m)$ such that
\[
  v\le Cu-(C-1)\phi_m\ \text{quasi-everywhere}
  \qquad\text{and}\qquad
  \int_V(\beta_m)_{v}^k\ge2^{-k-1}h .
\]
\end{lemma}

\begin{proof}
Put $a_m\coloneqq \int_V(\beta_m)_u^k$ and
$b_m\coloneqq \int_V(\beta_m)_{\phi_m}^k$. By multilinearity,
$a_m=\sum_{j=0}^k\binom kj\eta_m^{k-j}\int_V\gamma_u^j\wedge\omega_V^{k-j}\to h$,
and $b_m\to h$ by \cref{lem:trace-approximants}\ref{it:vq-mass}. Hence,
for $m$ large, $u,\phi_m\in\PSH(V,\beta_m)_{>0}$,
$(2C)^ka_m>((2C)^k-1)b_m$, $b_m>h/2$, and $u\lessP\phi_m$ by
\cref{lem:trace-approximants}\ref{it:vq-order}. \Cref{thm:rooftop},
applied on $V$ with $\theta=\beta_m$, $\varphi=u$, $\psi=\phi_m$ and
$t=2C$, shows that
$w\coloneqq P_{\beta_m}\bigl(2Cu-(2C-1)\phi_m\bigr)$ belongs to
$\PSH(V,\beta_m)$; we only use that it is not identically $-\infty$. Put
$v\coloneqq \tfrac12(w+\phi_m)\in\PSH(V,\beta_m)$. Since
$w\le 2Cu-(2C-1)\phi_m$ quasi-everywhere, $v\le Cu-(C-1)\phi_m$
quasi-everywhere. By multilinearity and the positivity of mixed
non-pluripolar products,
\[
  \int_V(\beta_m)_{v}^k
  =\int_V\Bigl(\tfrac12(\beta_m)_{w}+\tfrac12(\beta_m)_{\phi_m}\Bigr)^k
  \ge2^{-k}b_m>2^{-k-1}h. \qedhere
\]
\end{proof}

\subsection{The first level}\label{sec:first-level}

We now fix a dyadic integer $m$ such that \cref{lem:positive-rooftop}
applies and $\eta_m\le\delta_*/2$, and let $v$ be as in that lemma.
\Cref{lem:domination}, applied with $\eta=\eta_m\le1$, gives a closed
positive current $K^+\in\{\beta_m\}$ with $K^+\ge\delta_*\omega_V$ and
$K^+\lessI v$. Hence
\begin{equation}\label{eq:def-K}
  K\coloneqq K^+-\eta_m\omega_V=\gamma+\ddc\varphi_K\in\alpha|_V,\qquad
  K\ge4\delta\,\omega_V,\qquad
  \varphi_K\lessI v .
\end{equation}
Next, \cref{lem:ideal-function}, applied to the nonzero ideal
$\mathfrak j\mathfrak B_m$ on $X$, gives a quasi-psh function $\Psi_m$ on
$X$ and a constant $c_\Psi>0$ with
\begin{equation}\label{eq:def-Psi}
  \Psi_m=\psi_{\mathfrak j}+\psi_{\mathfrak B_m}+O(1)\ \text{locally},
  \qquad
  \ddc\Psi_m\ge-c_\Psi\,\omega .
\end{equation}
In particular, $\Psi_m$ is locally bounded outside
$Z'\cup Z(\mathfrak B_m)$, where $Z(\mathfrak B_m)$ is the zero locus of
$\mathfrak B_m$. Finally, since
$\varphi_K\in\PSH(V,\gamma-4\delta\omega_V)$, \cite[Theorem~1.6.2]{XiaSGPT}
provides a quasi-equisingular approximation $(\varphi_{K,q})_q$ of
$\varphi_K$, indexed by the dyadic integers $q$, such that
$\varphi_{K,q}$ has analytic singularities of type
$(q^{-1},\mathcal I(q\varphi_K))$ and
$\gamma+\ddc\varphi_{K,q}\ge(4\delta-\eta'_q)\omega_V$ with
$\eta'_q\downarrow0$.

\subsection{The second level}

Choose a dyadic integer $\ell\ge m$ such that $\eta'_\ell\le2\delta$ and
such that
\begin{equation}\label{eq:choice-M}
  M\coloneqq C\ell\quad\text{satisfies}\quad
  M>L\quad\text{and}\quad M\ge(C-1)c_\Psi/\delta .
\end{equation}
Since $m$ divides $\ell$, the numbers $M/m$ and $p\coloneqq (C-1)\ell/m$
are integers, and $M/m=p+\ell/m$. Put
\begin{equation}\label{eq:def-K-ell}
  K_\ell\coloneqq \gamma+\ddc\varphi_{K,\ell}\ge2\delta\,\omega_V,
  \qquad
  \mathfrak k\coloneqq \mathcal I(\ell\varphi_K)\subset\mathcal O_V ,
\end{equation}
so that $\varphi_{K,\ell}$ has analytic singularities of type
$(\ell^{-1},\mathfrak k)$.

\begin{remark}[Order of choices]\label{rem:order}
The data are chosen in the following order: $\delta_*$ and $\delta$; the
current $R$ and the constants $\delta_R$, $A_R$ and $c_0$; the integers $L$
and $C$; the level $m$, and with it $K$, $\Psi_m$, $c_\Psi$ and
$(\varphi_{K,q})_q$; the level $\ell$, and with it $M$ and $p$; finally,
in \cref{sec:extension}, a modification $\pi$ and a rational number
$\tau$. Every constant measuring positivity is fixed before $M$, and the
conditions imposed on $\ell$ and $M$ are lower bounds.
\end{remark}

\begin{lemma}\label{lem:ideal-comparison}
The ideals $\mathfrak k$, $\mathfrak b_m$ and $\mathfrak b_M$ on $V$ satisfy
\begin{equation}\label{eq:ideal-comparison}
  \mathfrak k\,\mathfrak b_m^{\,p}\subset\mathcal I(Mu)\subset\mathfrak b_M
  \subset\mathfrak b_m^{\,M/m}.
\end{equation}
Consequently, on any modification of $V$ on which $\mathfrak k$,
$\mathfrak b_m$ and $\mathfrak b_M$ are invertible, with divisors $G_K$,
$G_m$ and $G_M$ respectively, the divisors $G_M-(M/m)G_m$ and
$F\coloneqq G_K+pG_m-G_M$ are effective, and
\begin{equation}\label{eq:GK-decomposition}
  \frac1\ell\,G_K=\frac1M\,G_M+\frac{C-1}M\Bigl(G_M-\frac Mm\,G_m\Bigr)
     +\frac CM\,F .
\end{equation}
\end{lemma}

\begin{proof}
The last two inclusions in \eqref{eq:ideal-comparison} are contained in
\eqref{eq:trace-ideals}. For the first, let $f$ be a local section of
$\mathfrak k$ and $g$ a local section of $\mathfrak b_m^{\,p}$. By
\eqref{eq:def-K}, $f\in\mathcal I(\ell v)$, that is,
$|f|^2\mathrm{e}^{-\ell v}$ is locally integrable. By
\eqref{eq:phiq-singularities} and $mp=(C-1)\ell$,
$|g|^2\mathrm{e}^{-(C-1)\ell\phi_m}\le c_g\,\mathrm{e}^{p\psi_{\mathfrak b_m}-mp\,\phi_m}$
is locally bounded outside the zero locus of $\mathfrak b_m$, hence almost
everywhere. Since $\ell v\le Mu-(C-1)\ell\phi_m$ almost
everywhere by \cref{lem:positive-rooftop},
\[
  |fg|^2\mathrm{e}^{-Mu}\le|f|^2\mathrm{e}^{-\ell v}\cdot
  |g|^2\mathrm{e}^{-(C-1)\ell\phi_m}
  \qquad\text{almost everywhere},
\]
so $fg\in\mathcal I(Mu)$.

On a modification as in the statement, \eqref{eq:ideal-comparison}
translates into $G_K+pG_m\ge G_M\ge(M/m)G_m$, which is the effectivity
of the two divisors. Finally, substituting $G_K=G_M-pG_m+F$,
$M=C\ell$ and $p=(C-1)\ell/m$, the coefficients of $G_M$, $G_m$ and $F$
on both sides of \eqref{eq:GK-decomposition} are $C/M$,
$-(C-1)/m$ and $C/M$ respectively.
\end{proof}

\section{Extension across the submanifold}\label{sec:extension}

We keep the notation of \cref{sec:setup,sec:positive-current}. In this
section we construct a modification $\pi\colon Y\to X$, a point $y$ in the
strict transform $W$ of $V$, a neighborhood $U$ of $W$, and quasi-psh
functions $\varphi_G$ on $Y$ and $\varphi_U$ on $U$, to which
\cref{prop:gluing} applies with $\theta_Y\coloneqq \pi^*\theta$. This will
contradict $V\subset\EnK(\alpha)$.

\subsection{The ambient current}\label{sec:ambient}

Apply \cref{lem:regularization} to $V_\theta$ on $(X,\omega)$, with
$\zeta=\theta$ and the level $M$: we obtain a quasi-psh function $\Upsilon_M$ on
$X$ with
\[
  \theta+\ddc \Upsilon_M\ge-\frac{c_0}M\,\omega,\qquad
  \Upsilon_M\ge\frac1M\,\psi_{\mathfrak B_M}+O(1)\ \text{locally on }X.
\]
Consider the current
$Q\coloneqq (1-L/M)(\theta+\ddc \Upsilon_M)+(L/M)R\in\alpha$, with potential
$\varphi_Q\coloneqq (1-L/M)\Upsilon_M+(L/M)\Phi_R$. Since $M>L$ and
$L\delta_R>c_0$,
\begin{equation}\label{eq:QM-positivity}
  Q\ge\frac{L\delta_R-(1-L/M)c_0}M\,\omega\ge\epsilon_0\,\omega,
  \qquad \epsilon_0\coloneqq \frac{L\delta_R-c_0}M>0 .
\end{equation}
Since the polar set of $R$ contains $V$, $\varphi_Q=-\infty$ on $V$, and
by \eqref{eq:PhiR-lower},
\begin{equation}\label{eq:phiQ-lower}
  \varphi_Q\ge\frac{1-L/M}M\,\psi_{\mathfrak B_M}
   +\frac{LA_R}M\bigl(\psi_{\mathcal I_V}+\psi_{\mathfrak j}\bigr)+O(1)
   \qquad\text{locally on }X.
\end{equation}

\subsection{The modification}\label{sec:modification}

We construct $\pi$ in three stages
\[
  Y=Y_3\longrightarrow Y_2\longrightarrow Y_1\longrightarrow Y_0=X,
\]
each obtained from \cref{lem:relative-principalization} applied to the
strict transform $W_{i-1}$ of $V$ in $Y_{i-1}$; we write $W_i$ for the
strict transform of $V$ in $Y_i$, $W\coloneqq W_3$, and $H_i$ for the
divisor on $Y_i$ given by the lemma at the $i$-th stage. Divisors on
$Y_i$ and on $W_i$ are tacitly replaced by their total transforms on the
later stages; since no divisor below contains the relevant strict
transform of $V$, this is compatible with restriction to the strict
transforms.

\emph{Stage 1.} Apply \cref{lem:relative-principalization} to
$\mathfrak B_m$ on $X$, with $V$ as submanifold. We obtain $Y_1$, $W_1$,
$H_1$, and an effective divisor $\widetilde G_m\ge H_1$ with
$\mathfrak B_m\mathcal O_{Y_1}=\mathcal O_{Y_1}(-\widetilde G_m)$. Its
restriction $G_m\coloneqq \widetilde G_m|_{W_1}$ is the divisor of
$\mathfrak b_m\mathcal O_{W_1}$.

\emph{Stage 2.} Since $\mathfrak B_M\subset\mathfrak B_m^{M/m}$ by
\eqref{eq:trace-ideals} and $\mathfrak B_m^{M/m}\mathcal O_{Y_1}$ is
invertible, there is a coherent ideal $\mathfrak c\subset\mathcal O_{Y_1}$
with
$\mathfrak B_M\mathcal O_{Y_1}=\mathcal O_{Y_1}\bigl(-\tfrac Mm\widetilde G_m\bigr)\cdot\mathfrak c$.
Restricting to $W_1$ gives
$\mathfrak b_M\mathcal O_{W_1}=\mathcal O_{W_1}\bigl(-\tfrac MmG_m\bigr)\cdot\rest{W_1}\mathfrak c$,
so $\rest{W_1}\mathfrak c\ne0$. Apply
\cref{lem:relative-principalization} to $\mathfrak c$ on $Y_1$, with
$W_1$ as submanifold. We obtain $Y_2$, $W_2$ and $H_2$ with
$\mathfrak c\mathcal O_{Y_2}=\mathcal O_{Y_2}(-H_2-H_2')$. Thus
$\mathfrak B_M\mathcal O_{Y_2}=\mathcal O_{Y_2}(-\widetilde G_M)$ with
$\widetilde G_M\coloneqq \tfrac Mm\widetilde G_m+H_2+H_2'$, the restriction
$G_M\coloneqq \widetilde G_M|_{W_2}$ is the divisor of
$\mathfrak b_M\mathcal O_{W_2}$, and
\cref{lem:relative-principalization}\ref{it:rp-restriction} gives
\[
  H_2|_{W_2}=G_M-\frac Mm\,G_m .
\]

\emph{Stage 3.} On $W_2$ the ideals $\mathfrak b_m^{\,p}\mathcal O_{W_2}=\mathcal
O_{W_2}(-pG_m)$ and $\mathfrak b_M\mathcal O_{W_2}=\mathcal O_{W_2}(-G_M)$
are invertible, so \eqref{eq:ideal-comparison} shows that
\[
  \mathfrak f\coloneqq (\mathfrak k\mathcal O_{W_2})\cdot
  \mathcal O_{W_2}(G_M-pG_m)
\]
is a nonzero coherent ideal of $\mathcal O_{W_2}$. Let
$\widetilde{\mathfrak f}\subset\mathcal O_{Y_2}$ be its inverse image under
$\mathcal O_{Y_2}\to\mathcal O_{W_2}$, so that
$\rest{W_2}\widetilde{\mathfrak f}=\mathfrak f\ne0$. Apply
\cref{lem:relative-principalization} to $\widetilde{\mathfrak f}$ on
$Y_2$, with $W_2$ as submanifold. We obtain $Y=Y_3$, $W=W_3$ and $H_3$,
and \ref{it:rp-restriction} of that lemma shows that
$\mathfrak f\mathcal O_W=\mathcal O_W(-F)$ with $F\coloneqq H_3|_W$. Hence
$\mathfrak k\mathcal O_W=\mathcal O_W(-G_K)$ with $G_K=G_M-pG_m+F$, in
accordance with \cref{lem:ideal-comparison}.

Let $\pi\colon Y\to X$ be the composition. It is adapted to $V$, so
$\pi|_W\colon W\to V$ is a modification of $V$ and no $\pi$-exceptional
divisor contains $W$. By
\cref{lem:relative-principalization}\ref{it:rp-ideal-W}, applied to the
three stages,
\begin{equation}\label{eq:IV-H}
  \mathcal I_V\mathcal O_Y=\mathcal O_Y(-H)\cdot\mathcal I_W,
  \qquad H\coloneqq H_1+H_2+H_3 .
\end{equation}
To summarize, on $Y$ we have
$\mathfrak B_m\mathcal O_Y=\mathcal O_Y(-\widetilde G_m)$ with
$\widetilde G_m\ge H_1$ and
$\mathfrak B_M\mathcal O_Y=\mathcal O_Y(-\widetilde G_M)$, and on $W$
\begin{equation}\label{eq:restrictions}
  \widetilde G_M|_W=G_M,\qquad
  H_2|_W=G_M-\frac Mm\,G_m,\qquad
  H_3|_W=F,
\end{equation}
where $G_m$, $G_M$, $G_K$ are the divisors of $\mathfrak b_m\mathcal O_W$,
$\mathfrak b_M\mathcal O_W$, $\mathfrak k\mathcal O_W$, and $F$ is as in
\cref{lem:ideal-comparison}.

\subsection{The global potential}

By \cref{lem:exceptional-metric}, applied to $\pi$ and $\omega$, there is an
effective $\pi$-exceptional $\mathbf Q$-divisor $E_\pi$, with a smooth
metric, such that $\omega_Y\coloneqq \pi^*\omega-\Theta_{E_\pi}$ is a
K\"ahler form on $Y$. Fix a rational number $\tau$ with
$0<\tau<\min\{\epsilon_0,\delta\}$, and put
\[
  \varphi_G\coloneqq \varphi_Q\circ\pi+\tau\psi_{E_\pi}.
\]
By \eqref{eq:QM-positivity},
\[
  \theta_Y+\ddc\varphi_G
  =\pi^*Q+\tau\bigl([E_\pi]-\Theta_{E_\pi}\bigr)
  \ge(\epsilon_0-\tau)\pi^*\omega+\tau\omega_Y
  \ge\tau\omega_Y ,
\]
and $\varphi_G=-\infty$ on $W$, because $\varphi_Q=-\infty$ on $V$. Let
$\Lambda\subset Y\setminus W$ be compact. Pulling back
\eqref{eq:phiQ-lower}, and using
$\psi_{\mathfrak B_M}\circ\pi=\psi_{\widetilde G_M}+O(1)$ and
$\psi_{\mathcal I_V}\circ\pi=\psi_H+\psi_{\mathcal I_W}+O(1)$ (by
\eqref{eq:IV-H}), where $\psi_{\mathcal I_W}$ is bounded on $\Lambda$, we
obtain a constant $C_\Lambda$ with
\begin{equation}\label{eq:phiG-lower}
  \varphi_G\ge\frac{1-L/M}M\,\psi_{\widetilde G_M}
   +\frac{LA_R}M\bigl(\psi_H+\psi_{\mathfrak j}\circ\pi\bigr)
   +\tau\psi_{E_\pi}-C_\Lambda
  \qquad\text{on }\Lambda .
\end{equation}

\subsection{The local potential}

Define the effective $\mathbf Q$-divisor
\begin{equation}\label{eq:def-D}
  D\coloneqq \frac1M\,\widetilde G_M+\frac{C-1}M\,H_2+\frac CM\,H_3
  +\tau E_\pi
\end{equation}
on $Y$. No component of $D$ contains $W$, and by \eqref{eq:restrictions}
and \eqref{eq:GK-decomposition},
\begin{equation}\label{eq:D-restriction}
  D|_W=\frac1M\,G_M+\frac{C-1}M\Bigl(G_M-\frac Mm\,G_m\Bigr)+\frac CM\,F
       +\tau E_\pi|_W
  =\frac1\ell\,G_K+\tau E_\pi|_W .
\end{equation}
This identity is the point of the construction: the divisor $D$ on $Y$
restricts on $W$ to the singularity divisor of the pullback of $K_\ell$
(up to the small correction $\tau E_\pi$), while its coefficients along
$H_2$ and $H_3$ are large enough to dominate the singularities of
$\varphi_G$ in \eqref{eq:phiG-lower}.

\emph{A K\"ahler class on $W$.} Put
$\theta_W\coloneqq (\theta_Y-\delta\pi^*\omega)|_W$ and
$\varphi'\coloneqq (\pi|_W)^*\varphi_{K,\ell}+\tau\psi_{E_\pi}|_W$. Since
$\varphi_{K,\ell}$ has analytic singularities of type
$(\ell^{-1},\mathfrak k)$ and $\mathfrak k\mathcal O_W=\mathcal O_W(-G_K)$,
the function $\varphi'$ has log singularities along $D|_W$ by
\eqref{eq:D-restriction}. By \eqref{eq:def-K-ell} and $\tau<\delta$,
\[
  \theta_W+\ddc\varphi'
  =(\pi|_W)^*K_\ell-\delta\pi^*\omega|_W
   +\tau\bigl([E_\pi|_W]-\Theta_{E_\pi}|_W\bigr)
  \ge(\delta-\tau)\pi^*\omega|_W+\tau\omega_Y|_W
  \ge\tau\omega_Y|_W .
\]
By \cref{lem:residual-nef}, the class
$\{\theta_W\}-\{D|_W\}=\{\theta_Y-\Theta_D-\delta\pi^*\omega\}|_W$ is
K\"ahler on $W$.

\emph{Extension to a neighborhood of $W$.} Choose a K\"ahler form on $W$
in this class, of the form
$(\theta_Y-\Theta_D-\delta\pi^*\omega)|_W+\ddc f$ with $f$ smooth on $W$;
extend $f$ to a smooth function $\widetilde f$ on $Y$, and let $r_W^2$ be
the squared distance to $W$ for a Riemannian metric on $Y$, which is smooth
near $W$. At points of $W$, the form $\ddc r_W^2$ is semipositive with
kernel the tangent space of $W$. Hence, by compactness of $W$, for $b\gg1$
the form
\[
  \Omega\coloneqq \theta_Y-\Theta_D-\delta\pi^*\omega
  +\ddc\bigl(\widetilde f+b\,r_W^2\bigr)
\]
is positive at every point of $W$, and there are a relatively compact
open neighborhood $U$ of $W$, on which $r_W^2$ is smooth, and a constant
$c_\Omega>0$ with $\Omega\ge c_\Omega\omega_Y$ on $U$. Put
\[
  \chi\coloneqq \frac{C-1}M\,\Psi_m\circ\pi,
  \qquad
  \varphi_U\coloneqq \widetilde f+b\,r_W^2+\psi_D+\chi\quad\text{on }U.
\]
By \eqref{eq:def-Psi} and \eqref{eq:choice-M},
$\ddc\chi\ge-\frac{(C-1)c_\Psi}M\pi^*\omega\ge-\delta\pi^*\omega$, so
\[
  \theta_Y+\ddc\varphi_U
  =\Omega+[D]+\bigl(\delta\pi^*\omega+\ddc\chi\bigr)
  \ge c_\Omega\,\omega_Y
  \qquad\text{on }U.
\]

\emph{The point $y$.} Since no component of $D$ contains $W$, and since
$\pi(W)=V$ is not contained in $Z'\cup Z(\mathfrak B_m)$ (recall that
$\mathfrak B_m$ is the unit ideal near $x_*\in V$), the set of points of
$W$ lying in $\Supp D\cup\pi^{-1}\bigl(Z'\cup Z(\mathfrak B_m)\bigr)$ is a
proper analytic subset of $W$. Choose $y\in W$ outside it. Near $y$,
$\psi_D$ is smooth and $\chi$ is bounded, so $\varphi_U$ is bounded near
$y$.

\emph{The collar condition.} Let $\Lambda\subset U\setminus W$ be
compact. On $\Lambda$, the function $\widetilde f+b\,r_W^2$ is bounded.
By \eqref{eq:def-Psi} and $\widetilde G_m\ge H_1$, we have
$\chi\le\frac{C-1}M\bigl(\psi_{\mathfrak j}\circ\pi+\psi_{H_1}\bigr)+O(1)$,
and therefore, by \eqref{eq:def-D},
\[
  \psi_D+\chi\le\frac1M\,\psi_{\widetilde G_M}
  +\frac{C-1}M\bigl(\psi_{H_1}+\psi_{H_2}+\psi_{\mathfrak j}\circ\pi\bigr)
  +\frac CM\,\psi_{H_3}+\tau\psi_{E_\pi}+O(1).
\]
Compare this with \eqref{eq:phiG-lower}, in which
$\psi_H=\psi_{H_1}+\psi_{H_2}+\psi_{H_3}+O(1)$. The coefficients of
$\psi_{\widetilde G_M}$, of $\psi_{H_1}$, $\psi_{H_2}$ and
$\psi_{\mathfrak j}\circ\pi$, and of $\psi_{H_3}$ above are $1/M$, $(C-1)/M$
and $C/M$, and they are at least the corresponding coefficients
$(1-L/M)/M$, $LA_R/M$ and $LA_R/M$ in \eqref{eq:phiG-lower}, by
\eqref{eq:constants}. Since all these functions are locally bounded above,
it follows that $\psi_D+\chi\le\varphi_G+C'_\Lambda$, and hence
$\varphi_U\le\varphi_G+C''_\Lambda$ on $\Lambda$, for suitable constants.

\subsection{End of the proof}\label{sec:conclusion}

The manifold $Y$ is compact K\"ahler, $W$ is compact, and
\cref{prop:gluing} applies to $\theta_Y=\pi^*\theta$, $W$, $U$, $y$,
$\varphi_G$ and $\varphi_U$, with the K\"ahler form
$\omega_0\coloneqq \min\{\tau,c_\Omega\}\,\omega_Y$: hypotheses
\ref{it:gl-current}, \ref{it:gl-local} and \ref{it:gl-collar} were
verified in the last two subsections. Hence $y\notin\EnK(\pi^*\alpha)$. On
the other hand, by \eqref{eq:transformation-rules},
\[
  \EnK(\pi^*\alpha)\supset\pi^{-1}\bigl(\EnK(\alpha)\bigr)
  \supset\pi^{-1}(V)\supset W\ni y,
\]
a contradiction. Therefore the assumption $h>0$ made in
\cref{sec:setup} is false, and $\langle\alpha^k\rangle_{X\mid V}=0$. This
proves the first assertion of \cref{thm:main}.

For the second assertion, if $\alpha$ is not big, then $\EnK(\alpha)=X$ by
convention, and $\Null(\alpha)=X$ because
$\langle\alpha^n\rangle_{X\mid X}=\vol(\alpha)=0$. If $\alpha$ is big,
then $\EnK(\alpha)$ has no isolated points by
\cref{lem:no-isolated-points}, so each of its irreducible components is
positive-dimensional and, by the first assertion, contained in
$\Null(\alpha)$; thus $\EnK(\alpha)\subset\Null(\alpha)$. Conversely, a
positive-dimensional irreducible subvariety $V\not\subset\EnK(\alpha)$ has
positive numerical restricted volume by
\cref{sec:restricted-volumes}\ref{it:rv-positive}, so every subvariety
occurring in \eqref{eq:null-locus} is contained in $\EnK(\alpha)$, and
$\Null(\alpha)\subset\EnK(\alpha)$. This completes the proof of
\cref{thm:main}.\qed

\printbibliography

\bigskip
\noindent Duc Viet Vu, \textsc{Universit\"at zu K\"oln}\par\nopagebreak
\noindent\textit{Email address:} \texttt{vuvietsp@gmail.com}

\medskip
\noindent Mingchen Xia, \textsc{Institute of Geometry and Physics,
University of Science and Technology of China}\par\nopagebreak
\noindent\textit{Email address:} \texttt{xiamingchen2008@gmail.com}

\end{document}
